% EDITION STATUS: partial; source units 58/58
\documentclass[11pt,a4paper]{article}
\usepackage[no-math]{fontspec}
\setmainfont{Latin Modern Roman}
\setsansfont{Latin Modern Sans}
\setmonofont{Latin Modern Mono}
\usepackage[french,english]{babel}
\usepackage[left=22mm,right=34mm,top=22mm,bottom=22mm,marginparwidth=23mm,marginparsep=5mm,includeheadfoot]{geometry}
\usepackage{amsmath,amssymb}
\usepackage[normalem]{ulem}
\usepackage{xcolor,tikz-cd,marginnote,etoolbox,microtype,graphicx}
\usetikzlibrary{shapes.geometric}
\microtypesetup{protrusion=false}
\usepackage[unicode,colorlinks=true,linkcolor=LinkInk,urlcolor=LinkInk,bookmarksopen=true]{hyperref}
\definecolor{LinkInk}{HTML}{254E70}
\definecolor{HandInk}{HTML}{254E70}
\definecolor{DeletionInk}{HTML}{8B3A3A}
\definecolor{QuietInk}{HTML}{555555}
\setlength{\parindent}{1em}
\setlength{\parskip}{3pt plus 1pt}
\setlength{\emergencystretch}{1.5em}
\raggedbottom
\normalmarginpar
\frenchspacing
\renewcommand{\textcircled}[1]{%
  \tikz[baseline=(editionnumber.base)]{%
    \node[draw=.,text=.,circle,inner sep=.4pt,outer sep=0pt,
      line width=.35pt] (editionnumber) {#1};}}
\newif\ifdiplomatic
\newif\ifreading
\newif\ifenglish
\diplomaticfalse
\readingfalse
\englishtrue
\newcommand{\EditionTitle}{Sketch of a Program}
\newcommand{\EditionForm}{Corrected text with translator’s notes}
\hypersetup{pdftitle={Sketch of a Program — Corrected text with translator’s notes},pdfauthor={Alexander Grothendieck},pdfsubject={A scan-checked edition of Esquisse d'un programme, January 1984, Montpellier cote 119}}
\newcommand{\DeclareScan}[4]{%
  \expandafter\def\csname edurl@#1\endcsname{#2}%
  \expandafter\def\csname edkey@#1\endcsname{#3}%
  \expandafter\def\csname edpage@#1\endcsname{#4}}
\def\CurrentScanURL{}
\def\CurrentPageKey{}
\def\CurrentCopy{copy1}
\def\CurrentApparatus{}
\newcounter{appnote}
\newcommand{\EditionPart}[2]{\clearpage\def\CurrentCopy{#1}\section*{#2}}
\newcommand{\scanpage}[1]{%
  \clearpage
  \expandafter\let\expandafter\CurrentScanURL\csname edurl@#1\endcsname
  \expandafter\let\expandafter\CurrentPageKey\csname edkey@#1\endcsname
  \def\CurrentApparatus{}%
  \hypertarget{scan-\CurrentPageKey}{}%
  \typeout{EDITION-SCAN: \CurrentPageKey; PDFPAGE: \the\numexpr\ReadonlyShipoutCounter+1\relax}%
  \ifreading
    \marginnote{\raggedright\scriptsize\sffamily\href{\CurrentScanURL}{119 / \csname edpage@#1\endcsname}}%
  \else
    \noindent{\small\sffamily\href{\CurrentScanURL}{[scan #1]}}\par\smallskip
  \fi}
\newcommand{\mspage}[1]{%
  \ifreading\marginnote{\raggedright\scriptsize\sffamily p.\,#1}[2\baselineskip]%
  \else\noindent{\small\sffamily[p.\,#1]}\par\smallskip\fi}
% A hand macro around a tikz-cd arrow marks that edge only. TikZ defers edge
% drawing, so its explicit color option is required; a surrounding text color
% group alone would not survive until the arrow is drawn.
\ExplSyntaxOn
\NewDocumentCommand{\EditionHand}{m}
  {
    \tl_if_head_eq_meaning:nNTF {#1} \arrow
      {\ifdiplomatic \EditionColorArrow #1 \else #1\fi}
      {\ifdiplomatic \textcolor{HandInk}{#1}\else #1\fi}
  }
\cs_new_protected:Npn \EditionColorArrow \arrow [#1]
  {\arrow[#1,color=HandInk]}
\ExplSyntaxOff
\newcommand{\hand}[1]{\EditionHand{#1}}
% The source oval is handwriting; its already typed contents retain black ink.
\newcommand{\handcircle}[1]{%
  \tikz[baseline=(editioncircled.base)]{%
    \node[draw=HandInk,ellipse,inner xsep=.9pt,inner ysep=.2pt,
      outer sep=0pt,line width=.4pt,text=black] (editioncircled) {#1};}}

\let\over\relax
\newcommand{\over}[2]{\textcolor{QuietInk}{#1}\textcolor{HandInk}{\ensuremath{^{\!\nearrow}}#2}}
\newcommand{\inserted}[2][]{\textcolor{HandInk}{\ensuremath{\langle}#2\ensuremath{\rangle}}}
\newcommand{\deleted}[1]{\ifmmode\textcolor{DeletionInk}{[#1]_{\scriptscriptstyle\mathrm{del}}}\else\textcolor{DeletionInk}{\sout{#1}}\fi}
\renewcommand{\xout}[1]{\ifmmode\textcolor{QuietInk}{[#1]_{\scriptscriptstyle\mathrm{x}}}\else\textcolor{QuietInk}{\sout{#1}}\fi}
\let\TextUline\uline
\let\TextUuline\uuline
\makeatletter
\UL@pixel=0pt
\makeatother
\renewcommand{\uline}[1]{\ifmmode\underline{#1}\else\TextUline{#1}\fi}
\renewcommand{\uuline}[1]{\ifmmode\underline{\underline{#1}}\else\TextUuline{#1}\fi}
\newcommand{\unsure}[1]{\href{\CurrentScanURL}{[#1?]}}
\newcommand{\illegible}[1]{\href{\CurrentScanURL}{[\ifenglish illegible\else illisible\fi\ #1]}}
\newsavebox{\MarginalTestBox}
\newdimen\MarginalLineHeight
\newcount\MarginalLineCount
\newif\ifMarginalFits
\newif\ifMarginCluster
\newdimen\MarginalClusterOffset
\newdimen\MarginalPlacementOffset
\newcommand{\margincluster}[1]{%
  \begingroup\MarginClustertrue\global\MarginalClusterOffset=0pt #1\endgroup}
\newcommand{\margin}[1]{%
  \begingroup
  \let\MarginalSavedApparatus\CurrentApparatus
  \edef\MarginalSavedAppNumber{\the\value{appnote}}%
  \setbox\MarginalTestBox=\vbox{%
    \hsize=\marginparwidth\linewidth=\marginparwidth
    \parindent=0pt\parskip=0pt\footnotesize
    \global\MarginalLineHeight=\baselineskip
    \raggedright\noindent\textcolor{HandInk}{#1}\par}%
  \directlua{
    local lines = 0
    for n in node.traverse_id(node.id("hlist"), tex.box["MarginalTestBox"].head) do
      lines = lines + 1
    end
    tex.count["MarginalLineCount"] = lines
  }%
  \MarginalFitstrue
  \ifnum\MarginalLineCount>3\relax\MarginalFitsfalse\fi
  \ifdim\dimexpr\ht\MarginalTestBox+\dp\MarginalTestBox\relax>3\MarginalLineHeight\relax\MarginalFitsfalse\fi
  \ifMarginalFits
    \MarginalPlacementOffset=0pt
    \ifMarginCluster\MarginalPlacementOffset=\MarginalClusterOffset\fi
    \typeout{EDITION-MARGIN: \CurrentPageKey; PLACEMENT: margin; LINES: \the\MarginalLineCount; BOX-HEIGHT: \the\dimexpr\ht\MarginalTestBox+\dp\MarginalTestBox\relax; BASELINE: \the\MarginalLineHeight; LABELLED: no; OFFSET: \the\MarginalPlacementOffset}%
    \marginnote{\usebox{\MarginalTestBox}}[\MarginalPlacementOffset]%
    \ifMarginCluster\global\advance\MarginalClusterOffset by\dimexpr\ht\MarginalTestBox+\dp\MarginalTestBox+3pt\relax\fi
  \else
    \typeout{EDITION-MARGIN: \CurrentPageKey; PLACEMENT: block; LINES: \the\MarginalLineCount; BOX-HEIGHT: \the\dimexpr\ht\MarginalTestBox+\dp\MarginalTestBox\relax; BASELINE: \the\MarginalLineHeight; LABELLED: yes}%
    \setcounter{appnote}{\MarginalSavedAppNumber}%
    \global\let\CurrentApparatus\MarginalSavedApparatus
    \par\nobreak\smallskip\nobreak\noindent\begin{minipage}{\linewidth}\hrule\smallskip
    {\footnotesize\textcolor{HandInk}{\ifenglish Marginal note:\else Note marginale :\fi\ #1}}\smallskip\hrule\end{minipage}\par\smallskip
    \ifMarginCluster\global\MarginalClusterOffset=0pt\fi
  \fi
  \endgroup}
\newcommand{\authornote}[1]{\margin{#1}}
\newcommand{\diagram}[1]{\par\smallskip{\centering#1\par}\smallskip}
\newenvironment{diagramblock}{\par\begin{center}}{\end{center}\par}
\newcommand{\DeclareVariant}[2]{\expandafter\gdef\csname edvariant@#1\endcsname{#2}}
\newcommand{\AppItem}[2]{%
  \ifcsdef{edvariant@#2}{%
    \ifcsempty{edvariant@#2}{\PackageError{Esquisse}{Empty apparatus entry #2}{Every executed apparatus marker requires its printed note.}}{}%
  }{\PackageError{Esquisse}{Undefined apparatus entry #2}{Every executed apparatus marker requires its printed note.}}%
  \typeout{EDITION-APP-ENTRY: #2; NUMBER: #1}%
  \par\noindent\textsuperscript{#1}\csname edvariant@#2\endcsname}
\newcommand{\EditionApp}[1]{%
  \stepcounter{appnote}%
  \ifmmode{}^{\scriptscriptstyle\theappnote}\else\textsuperscript{\theappnote}\fi
  \begingroup\edef\x{\endgroup\noexpand\gappto\noexpand\CurrentApparatus{\noexpand\AppItem{\theappnote}{#1}}}\x}
\newcommand{\var}[4]{\ifdiplomatic#1\else#2\ifreading\else\EditionApp{#4}\fi\fi}
\newcommand{\scancontext}[2]{\begingroup\def\CurrentScanURL{#1}#2\endgroup}
\newcommand{\copytwo}[2]{\ifdiplomatic\else\ifreading\else\DeclareVariant{#1}{#2}\EditionApp{#1}\fi\fi}
\newcommand{\frapparatus}[1]{\ifreading\else\EditionApp{#1}\fi}
\newcommand{\tnote}[2]{\ifreading\else\DeclareVariant{#1}{\textit{Translator's note.} #2}\EditionApp{#1}\fi}
\newcommand{\finishscan}{\ifx\CurrentApparatus\empty\else\par\bigskip\noindent\rule{.35\linewidth}{.35pt}\par\smallskip{\footnotesize\setlength{\parskip}{3pt}\CurrentApparatus}\fi}
\newcommand{\DeclareNoteReturns}[2]{\expandafter\def\csname edreturns@#1\endcsname{#2}}
\newcommand{\notecall}[3]{\hypertarget{#2}{\hyperlink{note-#1}{#3}}}
\newcommand{\notedef}[2]{\hypertarget{note-#1}{#2}\,{\scriptsize[\csname edreturns@#1\endcsname]}}
\newcommand{\AppReading}[1]{{\def\margin##1{##1}\def\authornote##1{##1}#1}}
\newcommand{\EmptyReading}{\ifenglish\textit{omitted}\else\textit{supprimé}\fi}
% Observed rounded source glyph, cote 119, site 42, PDF page 43.
% Its conventional alphabet assignment remains uncertain; use with \unsure.
% The contour is a vector traced from the reviewed scan detail, not a font alias.
\makeatletter
\newcommand{\sourceglyph}[1]{\mathord{\mathpalette{\EditionSourceGlyph{#1}}{}}}
\newcommand{\EditionSourceGlyph}[3]{%
  \begingroup
  \setbox0=\hbox{$#2 A$}%
  \resizebox{!}{\ht0}{\csname EditionGlyph@#1\endcsname}%
  \endgroup}
\expandafter\newcommand\csname EditionGlyph@119A\endcsname{%
  \begin{tikzpicture}[x=1pt,y=1pt,baseline=0pt]
  \path[fill,even odd rule]
(3.450,10.100) -- (4.950,10.100) -- (5.050,10.000) -- (5.550,10.000) -- (5.650,9.900) -- (5.950,9.900) -- (6.050,9.800) -- (6.250,9.800) -- (6.650,9.600) -- (6.950,9.300) -- (7.050,9.300) -- (7.600,8.750) -- (7.800,8.350) -- (7.800,8.050) -- (7.550,7.800) -- (7.050,7.800) -- (6.250,8.400) -- (6.050,8.500) -- (5.850,8.500) -- (5.350,8.800) -- (3.950,8.800) -- (3.850,8.900) -- (3.650,8.900) -- (3.550,8.800) -- (3.350,8.800) -- (3.150,8.700) -- (2.250,7.900) -- (2.100,7.850) -- (2.100,7.750) -- (1.900,7.550) -- (1.500,6.750) -- (1.500,6.450) -- (1.400,6.350) -- (1.400,5.850) -- (1.300,5.750) -- (1.300,4.650) -- (1.200,4.550) -- (1.200,4.050) -- (1.300,3.950) -- (1.300,2.450) -- (1.400,2.350) -- (1.400,2.150) -- (1.700,1.850) -- (1.900,1.450) -- (1.900,1.250) -- (2.000,1.050) -- (2.150,1.000) -- (2.550,0.600) -- (2.850,0.500) -- (3.100,0.250) -- (3.100,0.150) -- (2.850,0.000) -- (2.250,0.000) -- (2.050,0.100) -- (1.750,0.400) -- (1.650,0.400) -- (0.900,1.150) -- (0.900,1.350) -- (0.700,1.650) -- (0.700,1.850) -- (0.500,2.050) -- (0.300,2.450) -- (0.300,2.750) -- (0.200,2.850) -- (0.200,3.150) -- (0.100,3.250) -- (0.100,3.850) -- (0.000,3.950) -- (0.000,6.650) -- (0.300,7.050) -- (0.400,7.250) -- (0.400,7.450) -- (0.500,7.550) -- (0.500,7.750) -- (0.700,7.950) -- (0.800,8.250) -- (1.950,9.300) -- (2.050,9.300) -- (2.550,9.700) -- (3.150,10.000) -- (3.350,10.000) -- (3.450,10.100) -- cycle
(6.050,5.700) -- (6.250,5.700) -- (6.500,5.450) -- (6.500,5.250) -- (6.600,5.150) -- (6.600,4.650) -- (6.500,4.550) -- (6.500,2.850) -- (6.400,2.750) -- (6.500,2.550) -- (6.500,2.050) -- (6.600,1.950) -- (6.600,1.250) -- (6.950,0.900) -- (7.500,0.750) -- (7.400,0.550) -- (7.400,0.350) -- (7.250,0.100) -- (6.850,0.100) -- (6.250,0.400) -- (5.850,0.400) -- (5.150,0.000) -- (4.750,0.000) -- (4.650,0.100) -- (4.550,0.000) -- (3.650,0.000) -- (3.500,0.150) -- (3.550,0.300) -- (3.750,0.400) -- (4.250,0.400) -- (4.350,0.500) -- (4.850,0.500) -- (4.950,0.600) -- (5.150,0.600) -- (5.500,0.950) -- (5.600,1.150) -- (5.600,1.350) -- (5.700,1.450) -- (5.700,1.950) -- (5.600,2.050) -- (5.600,2.350) -- (5.500,2.450) -- (5.500,2.750) -- (5.400,2.850) -- (5.400,3.350) -- (5.500,3.450) -- (5.500,4.050) -- (5.600,4.150) -- (5.600,5.050) -- (5.800,5.450) -- (6.050,5.700) -- cycle
(7.950,1.900) -- (8.300,1.850) -- (8.300,1.050) -- (8.150,0.900) -- (7.950,0.900) -- (7.850,0.800) -- (7.550,0.800) -- (7.500,0.950) -- (7.600,1.050) -- (7.600,1.450) -- (7.950,1.900) -- cycle;
  \end{tikzpicture}}
\makeatother

\DeclareScan{119/2}{https://grothendieck.umontpellier.fr/web/viewer.html?file=../119.pdf\#page=3}{esq_119_002}{2}
\DeclareScan{119/3}{https://grothendieck.umontpellier.fr/web/viewer.html?file=../119.pdf\#page=4}{esq_119_003}{3}
\DeclareScan{119/4}{https://grothendieck.umontpellier.fr/web/viewer.html?file=../119.pdf\#page=5}{esq_119_004}{4}
\DeclareScan{119/5}{https://grothendieck.umontpellier.fr/web/viewer.html?file=../119.pdf\#page=6}{esq_119_005}{5}
\DeclareScan{119/6}{https://grothendieck.umontpellier.fr/web/viewer.html?file=../119.pdf\#page=7}{esq_119_006}{6}
\DeclareScan{119/7}{https://grothendieck.umontpellier.fr/web/viewer.html?file=../119.pdf\#page=8}{esq_119_007}{7}
\DeclareScan{119/8}{https://grothendieck.umontpellier.fr/web/viewer.html?file=../119.pdf\#page=9}{esq_119_008}{8}
\DeclareScan{119/9}{https://grothendieck.umontpellier.fr/web/viewer.html?file=../119.pdf\#page=10}{esq_119_009}{9}
\DeclareScan{119/10}{https://grothendieck.umontpellier.fr/web/viewer.html?file=../119.pdf\#page=11}{esq_119_010}{10}
\DeclareScan{119/11}{https://grothendieck.umontpellier.fr/web/viewer.html?file=../119.pdf\#page=12}{esq_119_011}{11}
\DeclareScan{119/12}{https://grothendieck.umontpellier.fr/web/viewer.html?file=../119.pdf\#page=13}{esq_119_012}{12}
\DeclareScan{119/13}{https://grothendieck.umontpellier.fr/web/viewer.html?file=../119.pdf\#page=14}{esq_119_013}{13}
\DeclareScan{119/14}{https://grothendieck.umontpellier.fr/web/viewer.html?file=../119.pdf\#page=15}{esq_119_014}{14}
\DeclareScan{119/15}{https://grothendieck.umontpellier.fr/web/viewer.html?file=../119.pdf\#page=16}{esq_119_015}{15}
\DeclareScan{119/16}{https://grothendieck.umontpellier.fr/web/viewer.html?file=../119.pdf\#page=17}{esq_119_016}{16}
\DeclareScan{119/17}{https://grothendieck.umontpellier.fr/web/viewer.html?file=../119.pdf\#page=18}{esq_119_017}{17}
\DeclareScan{119/18}{https://grothendieck.umontpellier.fr/web/viewer.html?file=../119.pdf\#page=19}{esq_119_018}{18}
\DeclareScan{119/19}{https://grothendieck.umontpellier.fr/web/viewer.html?file=../119.pdf\#page=20}{esq_119_019}{19}
\DeclareScan{119/20}{https://grothendieck.umontpellier.fr/web/viewer.html?file=../119.pdf\#page=21}{esq_119_020}{20}
\DeclareScan{119/21}{https://grothendieck.umontpellier.fr/web/viewer.html?file=../119.pdf\#page=22}{esq_119_021}{21}
\DeclareScan{119/22}{https://grothendieck.umontpellier.fr/web/viewer.html?file=../119.pdf\#page=23}{esq_119_022}{22}
\DeclareScan{119/23}{https://grothendieck.umontpellier.fr/web/viewer.html?file=../119.pdf\#page=24}{esq_119_023}{23}
\DeclareScan{119/24}{https://grothendieck.umontpellier.fr/web/viewer.html?file=../119.pdf\#page=25}{esq_119_024}{24}
\DeclareScan{119/25}{https://grothendieck.umontpellier.fr/web/viewer.html?file=../119.pdf\#page=26}{esq_119_025}{25}
\DeclareScan{119/26}{https://grothendieck.umontpellier.fr/web/viewer.html?file=../119.pdf\#page=27}{esq_119_026}{26}
\DeclareScan{119/27}{https://grothendieck.umontpellier.fr/web/viewer.html?file=../119.pdf\#page=28}{esq_119_027}{27}
\DeclareScan{119/28}{https://grothendieck.umontpellier.fr/web/viewer.html?file=../119.pdf\#page=29}{esq_119_028}{28}
\DeclareScan{119/29}{https://grothendieck.umontpellier.fr/web/viewer.html?file=../119.pdf\#page=30}{esq_119_029}{29}
\DeclareScan{119/30}{https://grothendieck.umontpellier.fr/web/viewer.html?file=../119.pdf\#page=31}{esq_119_030}{30}
\DeclareScan{119/31}{https://grothendieck.umontpellier.fr/web/viewer.html?file=../119.pdf\#page=32}{esq_119_031}{31}
\DeclareScan{119/32}{https://grothendieck.umontpellier.fr/web/viewer.html?file=../119.pdf\#page=33}{esq_119_032}{32}
\DeclareScan{119/33}{https://grothendieck.umontpellier.fr/web/viewer.html?file=../119.pdf\#page=34}{esq_119_033}{33}
\DeclareScan{119/34}{https://grothendieck.umontpellier.fr/web/viewer.html?file=../119.pdf\#page=35}{esq_119_034}{34}
\DeclareScan{119/35}{https://grothendieck.umontpellier.fr/web/viewer.html?file=../119.pdf\#page=36}{esq_119_035}{35}
\DeclareScan{119/36}{https://grothendieck.umontpellier.fr/web/viewer.html?file=../119.pdf\#page=37}{esq_119_036}{36}
\DeclareScan{119/37}{https://grothendieck.umontpellier.fr/web/viewer.html?file=../119.pdf\#page=38}{esq_119_037}{37}
\DeclareScan{119/38}{https://grothendieck.umontpellier.fr/web/viewer.html?file=../119.pdf\#page=39}{esq_119_038}{38}
\DeclareScan{119/39}{https://grothendieck.umontpellier.fr/web/viewer.html?file=../119.pdf\#page=40}{esq_119_039}{39}
\DeclareScan{119/40}{https://grothendieck.umontpellier.fr/web/viewer.html?file=../119.pdf\#page=41}{esq_119_040}{40}
\DeclareScan{119/41}{https://grothendieck.umontpellier.fr/web/viewer.html?file=../119.pdf\#page=42}{esq_119_041}{41}
\DeclareScan{119/42}{https://grothendieck.umontpellier.fr/web/viewer.html?file=../119.pdf\#page=43}{esq_119_042}{42}
\DeclareScan{119/43}{https://grothendieck.umontpellier.fr/web/viewer.html?file=../119.pdf\#page=44}{esq_119_043}{43}
\DeclareScan{119/44}{https://grothendieck.umontpellier.fr/web/viewer.html?file=../119.pdf\#page=45}{esq_119_044}{44}
\DeclareScan{119/45}{https://grothendieck.umontpellier.fr/web/viewer.html?file=../119.pdf\#page=46}{esq_119_045}{45}
\DeclareScan{119/46}{https://grothendieck.umontpellier.fr/web/viewer.html?file=../119.pdf\#page=47}{esq_119_046}{46}
\DeclareScan{119/47}{https://grothendieck.umontpellier.fr/web/viewer.html?file=../119.pdf\#page=48}{esq_119_047}{47}
\DeclareScan{119/48}{https://grothendieck.umontpellier.fr/web/viewer.html?file=../119.pdf\#page=49}{esq_119_048}{48}
\DeclareScan{119/49}{https://grothendieck.umontpellier.fr/web/viewer.html?file=../119.pdf\#page=50}{esq_119_049}{49}
\DeclareScan{119/50}{https://grothendieck.umontpellier.fr/web/viewer.html?file=../119.pdf\#page=51}{esq_119_050}{50}
\DeclareScan{119/51}{https://grothendieck.umontpellier.fr/web/viewer.html?file=../119.pdf\#page=52}{esq_119_051}{51}
\DeclareScan{119/52}{https://grothendieck.umontpellier.fr/web/viewer.html?file=../119.pdf\#page=53}{esq_119_052}{52}
\DeclareScan{119/53}{https://grothendieck.umontpellier.fr/web/viewer.html?file=../119.pdf\#page=54}{esq_119_053}{53}
\DeclareScan{119/54}{https://grothendieck.umontpellier.fr/web/viewer.html?file=../119.pdf\#page=55}{esq_119_054}{54}
\DeclareScan{119/55}{https://grothendieck.umontpellier.fr/web/viewer.html?file=../119.pdf\#page=56}{esq_119_055}{55}
\DeclareScan{119/56}{https://grothendieck.umontpellier.fr/web/viewer.html?file=../119.pdf\#page=57}{esq_119_056}{56}
\DeclareScan{119/57}{https://grothendieck.umontpellier.fr/web/viewer.html?file=../119.pdf\#page=58}{esq_119_057}{57}
\DeclareScan{119/58}{https://grothendieck.umontpellier.fr/web/viewer.html?file=../119.pdf\#page=59}{esq_119_058}{58}
\DeclareScan{119/59}{https://grothendieck.umontpellier.fr/web/viewer.html?file=../119.pdf\#page=60}{esq_119_059}{59}
\DeclareVariant{V-FR-010-001}{\textit{slip corrected}. \AppReading{Ladegraillerie} \(\rightarrow\) \AppReading{Ladegaillerie}. The preceding line gives Yves Ladegaillerie for the same participant; the second occurrence interposes an r.}
\DeclareVariant{V-FR-026-002}{\textit{slip corrected}. \AppReading{hyperquadratic} \(\rightarrow\) \AppReading{hyperquadric}. Copy 2 has a handwritten deletion of the two positions between hyperquadr and ique; the covered letters remain partly unreadable, compatible with at. Copy 1 clearly reads hyperquadratique. The body of the same page already names the circumscribed hyperquadrique. This deletion now supplies an explicit correction of the shared typing slip: hyperquadrique. See \href{https://grothendieck.umontpellier.fr/web/viewer.html?file=../119.pdf\#page=103}{copy 2, p.~24}.}
\DeclareVariant{V-FR-004-001}{\textit{French source apparatus: slip corrected}. Page reading: \AppReading{inadéquante}; corrected French: \AppReading{inadéquate}. Copy 2 has a handwritten slash through the n of inadéquante, a typing slip shared by both copies. Its correction gives inadéquate and enters the corrected copy1 text here. See \href{https://grothendieck.umontpellier.fr/web/viewer.html?file=../119.pdf\#page=81}{copy 2, p.~2, site 80}.}
\DeclareVariant{V-FR-006-001}{\textit{French source apparatus: slip corrected}. Page reading: \AppReading{des surface.}; corrected French: \AppReading{des surfaces.}. The local article des requires the plural; the same noun is written surfaces on page 3. The correction supplies the missing final s.}
\DeclareVariant{V-FR-006-002}{\textit{French source apparatus: slip corrected}. Page reading: \AppReading{d'automorphisme d'un groupe profini}; corrected French: \AppReading{d'automorphismes d'un groupe profini}. The text speaks of a group whose elements are automorphisms. Page 7 has the construction un groupe d'automorphismes; the final s is absent here in the same phrase.}
\DeclareVariant{V-FR-007-001}{\textit{French source apparatus: slip corrected}. Page reading: \AppReading{voir indispensable}; corrected French: \AppReading{voire indispensable}. The progression from more elegant to indispensable uses the intensifying adverb; voire occurs on page 2 (voire conformes) and page 3. The correction supplies its omitted final e here.}
\DeclareVariant{V-FR-009-001}{\textit{French source apparatus: slip corrected}. Page reading: \AppReading{convenanble}; corrected French: \AppReading{convenable}. An extra n is interposed immediately before b; convenable occurs in the same discussion on page 5. Copy 2 has a blue slash through that n of the typing slip shared by both copies, confirming the correction. See \href{https://grothendieck.umontpellier.fr/web/viewer.html?file=../119.pdf\#page=86}{copy 2, p.~7, site 85}.}
\DeclareVariant{V-FR-010-002}{\textit{French source apparatus: slip corrected}. Page reading: \AppReading{se trouve muni d'une décomposition}; corrected French: \AppReading{se trouve munie d'une décomposition}. The subject is chacune des douze pièces; the same page also uses munies for the surfaces. The exact feminine singular munie occurs on page 42. The final e is missing in this local agreement.}
\DeclareVariant{V-FR-012-001}{\textit{French source apparatus: slip corrected}. Page reading: \AppReading{Iladmet}; corrected French: \AppReading{Il admet}. The pronoun Il and verb admet are run together in the typing. The same page separately writes Il va and admet une description; the cartographic group is the subject of this sentence.}
\DeclareVariant{V-FR-013-001}{\textit{French source apparatus: slip corrected}. Page reading: \AppReading{les trois point 0, 1}; corrected French: \AppReading{les trois points 0, 1}. The number trois requires the plural; the same page earlier has de trois points, numérotés. The final s is missing in the second occurrence.}
\DeclareVariant{V-FR-014-001}{\textit{French source apparatus: slip corrected}. Page reading: \AppReading{ramifié seulement on}; corrected French: \AppReading{ramifié seulement en}. The preposition en introduces the ramification points; page 11 has the same construction, ramifiés au plus en les points. The o in on replaces e here.}
\DeclareVariant{V-FR-015-001}{\textit{French source apparatus: slip corrected}. Page reading: \AppReading{quels exactement}; corrected French: \AppReading{quelles exactement}. The explicit subject is the feminine plural les conjuguées de telle carte orientée; the next page has quelles sont les courbes. The proposal agrees the pronoun with that written subject.}
\DeclareVariant{V-FR-016-001}{\textit{French source apparatus: slip corrected}. Page reading: \AppReading{une directions}; corrected French: \AppReading{une direction}. The article une and adjective conceptuelle require the singular; page 25 has dans une direction toute différente. The s is added in error here.}
\DeclareVariant{V-FR-017-001}{\textit{French source apparatus: slip corrected}. Page reading: \AppReading{dingueen effet}; corrected French: \AppReading{dingue en effet}. The adjective dingue qualifies la supposition; en effet follows as an adverbial phrase. Page 19 explicitly separates en effet. A space is missing after dingue.}
\DeclareVariant{V-FR-019-001}{\textit{French source apparatus: slip corrected}. Page reading: \AppReading{est connu à isomorphisme près}; corrected French: \AppReading{est connue à isomorphisme près}. The written subject is une courbe algébrique anabélienne; the feminine form connue occurs on page 11. The final e is missing in this local agreement.}
\DeclareVariant{V-FR-024-001}{\textit{French source apparatus: slip corrected}. Page reading: \AppReading{nouveau ;}; corrected French: \AppReading{nouveaux ;}. Des charmes is plural; the form nouveaux occurs on page 20. The final x is missing in this occurrence.}
\DeclareVariant{V-FR-024-002}{\textit{French source apparatus: slip corrected}. Page reading: \AppReading{l'isosaèdre}; corrected French: \AppReading{l'icosaèdre}. The immediately following text gives un icosaèdre for the same polyhedron. The first occurrence has s in place of c.}
\DeclareVariant{V-FR-025-001}{\textit{French source apparatus: slip corrected}. Page reading: \AppReading{les caractéristique 2 et 5}; corrected French: \AppReading{les caractéristiques 2 et 5}. The article les and the two listed characteristics require the plural; the preceding line already has l'une des caractéristiques. The final s is absent here.}
\DeclareVariant{V-FR-026-001}{\textit{French source apparatus: slip corrected}. Page reading: \AppReading{réunisse même à le dégager}; corrected French: \AppReading{réussisse même à le dégager}. The context concerns succeeding in drawing out a principle, with à followed by an infinitive; page 46 writes je réussisse à communiquer in the same construction. The n in réunisse occupies the place of the double s in réussisse here.}
\DeclareVariant{V-FR-027-002}{\textit{French source apparatus: slip corrected}. Page reading: \AppReading{sphèrique}; corrected French: \AppReading{sphérique}. Page 13 gives sphérique with an acute accent, and the plural sphériques also occurs on page 26. The isolated grave accent in this occurrence is proposed as a typing slip.}
\DeclareVariant{V-FR-027-001}{\textit{French source apparatus: slip corrected}. Page reading: \AppReading{variété algèbrique}; corrected French: \AppReading{variété algébrique}. This page twice has géométrie algébrique with an acute accent, and also courbes algébriques. The isolated grave accent in variété algèbrique appears to be a typing slip; this proposal concerns that occurrence alone.}
\DeclareVariant{V-FR-028-001}{\textit{French source apparatus: slip corrected}. Page reading: \AppReading{rigidifidation}; corrected French: \AppReading{rigidification}. The term rigidification occurs on page 14 for the same kind of additional data; the page 26 occurrence has d in place of c.}
\DeclareVariant{V-FR-028-002}{\textit{French source apparatus: slip corrected}. Page reading: \AppReading{ces multiplicité}; corrected French: \AppReading{ces multiplicités}. Ces refers to the multiplicities named in the plural at the beginning of the page; the final s is missing in this later occurrence.}
\DeclareVariant{V-FR-028-003}{\textit{French source apparatus: slip corrected}. Page reading: \AppReading{s'identifisant au}; corrected French: \AppReading{s'identifiant au}. Copy 1 types s'identifisant; copy 2 has a handwritten stroke deleting precisely the s between identifi and ant, with both loops still visible. The deletion corrects that extra letter shared by the two typed copies, giving s'identifiant. See \href{https://grothendieck.umontpellier.fr/web/viewer.html?file=../119.pdf\#page=105}{copy 2, p.~26}.}
\DeclareVariant{V-FR-029-001}{\textit{French source apparatus: slip corrected}. Page reading: \AppReading{des homéomorphisme de telles paires}; corrected French: \AppReading{des homéomorphismes de telles paires}. The passage describes the morphisms in the category of pairs and uses plural des. Page 53 writes homéomorphismes; the final s is missing here. The later parallel supersedes the earlier editorial reservation.}
\DeclareVariant{V-FR-032-001}{\textit{French source apparatus: slip corrected}. Page reading: \AppReading{espaces semianalytique}; corrected French: \AppReading{espaces semianalytiques}. The adjective qualifies plural espaces; the next page gives semianalytiques for these same Hironaka spaces. The final s is missing here.}
\DeclareVariant{V-FR-033-001}{\textit{French source apparatus: slip corrected}. Page reading: \AppReading{stricte de toutes, celles des}; corrected French: \AppReading{stricte de toutes, celle des}. The sentence moves from the strictest theory to the broadest; vers celle qui and savoir celle on the same page supply the singular reference. The s in celles is proposed as an isolated addition.}
\DeclareVariant{V-FR-033-002}{\textit{French source apparatus: slip corrected}. Page reading: \AppReading{les catégories ``$\infty$-isotopique''}; corrected French: \AppReading{les catégories ``$\infty$-isotopiques''}. The same adjective is written in the plural earlier for the same categories; the final s is missing in this later occurrence.}
\DeclareVariant{V-FR-034-001}{\textit{French source apparatus: slip corrected}. Page reading: \AppReading{simple-%
mement}; corrected French: \AppReading{simple-%
ment}. The continuation mement repeats me at the line break; simplement is written on page 23. The correction removes that repetition while retaining the line break.}
\DeclareVariant{V-FR-035-001}{\textit{French source apparatus: slip corrected}. Page reading: \AppReading{mais à part}; corrected French: \AppReading{mis à part}. The passage sets the comparison theorems aside within a list; page 24 has Mis à part in the same expression. An extra a is interposed in mis here.}
\DeclareVariant{V-FR-035-002}{\textit{French source apparatus: slip corrected}. Page reading: \AppReading{toute application modéré}; corrected French: \AppReading{toute application modérée}. Toute application is feminine; page 37 gives une application modérée propre. The final e is missing in this occurrence.}
\DeclareVariant{V-FR-036-002}{\textit{French source apparatus: slip corrected}. Page reading: \AppReading{voisinages tubulaire}; corrected French: \AppReading{voisinages tubulaires}. The noun voisinages is plural; page 36 has the same phrase voisinages tubulaires. The final s is missing here.}
\DeclareVariant{V-FR-036-001}{\textit{French source apparatus: slip corrected}. Page reading: \AppReading{de façon toute aussi canonique}; corrected French: \AppReading{de façon tout aussi canonique}. Tout modifies aussi canonique here; page 8 already has tout aussi concrètes with a feminine adjective. The final e in toute is proposed as an isolated addition.}
\DeclareVariant{V-FR-038-001}{\textit{French source apparatus: slip corrected}. Page reading: \AppReading{j'étais parvenue}; corrected French: \AppReading{j'étais parvenu}. The same page gives auquel je suis parvenu for the same speaker and formalism; the e in parvenue is proposed as an isolated addition.}
\DeclareVariant{V-FR-039-001}{\textit{French source apparatus: slip corrected}. Page reading: \AppReading{d'espace modérés}; corrected French: \AppReading{d'espaces modérés}. The adjective modérés is plural and the map connects two spaces; the same passage writes deux espaces modérés. The noun lacks s in this occurrence.}
\DeclareVariant{V-FR-040-001}{\textit{French source apparatus: slip corrected}. Page reading: \AppReading{schématique (qui existent}; corrected French: \AppReading{schématiques (qui existent}. The adjective qualifies plural compactifications, followed by qui existent. Page 53 supplies the form schématiques. The correction adds only the missing final s in this occurrence.}
\DeclareVariant{V-FR-042-001}{\textit{French source apparatus: slip corrected}. Page reading: \AppReading{s'est agit de dégager}; corrected French: \AppReading{s'est agi de dégager}. Page 28 has exactly the construction il s'est agi de followed by an infinitive. The final t in agit is proposed as an isolated addition in the same construction on page 40; this parallel supersedes the earlier provisional reservation.}
\DeclareVariant{V-FR-043-001}{\textit{French source apparatus: slip corrected}. Page reading: \AppReading{divers expaces}; corrected French: \AppReading{divers espaces}. The same noun is written espaces earlier on this page. The occurrence expaces has x in place of s; the correction concerns that single letter.}
\DeclareVariant{V-FR-044-001}{\textit{French source apparatus: slip corrected}. Page reading: \AppReading{tout $\unsure{\sourceglyph{119A}}$-variété}; corrected French: \AppReading{toute $\unsure{\sourceglyph{119A}}$-variété}. Tout here determines the feminine noun variété. Page 14 has toute before the feminine noun courbe; the final e is missing in this occurrence. The distinctive glyph and its uncertainty are preserved.}
\DeclareVariant{V-FR-045-001}{\textit{French source apparatus: slip corrected}. Page reading: \AppReading{le manuscript}; corrected French: \AppReading{le manuscrit}. Page 45 gives manuscrit without p for another written work. The isolated form manuscript on page 43 has an extra p; only this occurrence is proposed for review, without changing other lexical choices.}
\DeclareVariant{V-FR-046-001}{\textit{French source apparatus: slip corrected}. Page reading: \AppReading{d'espacesd'Eilenberg}; corrected French: \AppReading{d'espaces d'Eilenberg}. Espaces and the elided preposition d' are run together; page 39 already separates espaces d'isomorphismes. The correction supplies only the missing space, preserving the following compound name.}
\DeclareVariant{V-FR-050-001}{\textit{French source apparatus: slip corrected}. Page reading: \AppReading{pen-%
dans ces dix}; corrected French: \AppReading{pen-%
dant ces dix}. The duration construction matches pendant une quinzaine on page 45. Only final s changes to t; the line break and the observed double space after dix remain intact.}
\DeclareVariant{V-FR-052-001}{\textit{French source apparatus: slip corrected}. Page reading: \AppReading{plus qu'en seul travail}; corrected French: \AppReading{plus qu'un seul travail}. The same paragraph counts written works and already has à un travail mieux. Only e in en changes to u to supply that article; the observed double space after goût remains intact.}
\DeclareVariant{V-FR-054-001}{\textit{French source apparatus: slip corrected}. Page reading: \AppReading{échappée jusqu'à présent}; corrected French: \AppReading{échappé jusqu'à présent}. The property is the subject of avait échappé and m' names the person it escaped. Page 34 has the same construction qui m'auraient échappé. The correction removes the isolated final e without changing subject or tense.}
\DeclareVariant{V-FR-055-001}{\textit{French source apparatus: slip corrected}. Page reading: \AppReading{(voir
même}; corrected French: \AppReading{(voire
même}. The passage extends the base from Q to Z using the intensifying adverb. The same phrase voire même occurs on pages 28, 32 and 34. The correction adds the missing e and preserves the line division; no formula changes.}
\DeclareVariant{V-FR-059-001}{\textit{French source apparatus: slip corrected}. Page reading: \AppReading{de l'espaces mo-%
}; corrected French: \AppReading{de l'espace mo-%
}. The article l' and adjective modéré refer to the single space named $X_*$. Page 37 writes l'espace in the same noun phrase. The added s is removed; the mo-déré line division and the formula remain unchanged.}
\DeclareVariant{V-FR-059-002}{\textit{French source apparatus: slip corrected}. Page reading: \AppReading{qui est munie d'une}; corrected French: \AppReading{qui est muni d'une}. The relative qui refers to the space $X_*$. Page 33 uses muni for the space T, and page 10 uses muni for K. The isolated final e is removed in this local agreement.}
\DeclareVariant{V-FR-059-003}{\textit{French source apparatus: slip corrected}. Page reading: \AppReading{est une plongement}; corrected French: \AppReading{est un plongement}. The preceding page has exactly un plongement. The noun and its adjective bordant are masculine in this occurrence; the extra e in une is removed.}
\DeclareNoteReturns{copy1-1}{\hyperlink{call-esq_119_007-1-1}{p.\,5}, \hyperlink{call-esq_119_044-1-2}{p.\,42}}
\DeclareNoteReturns{copy1-2}{\hyperlink{call-esq_119_014-2-1}{p.\,12}}
\DeclareNoteReturns{copy1-3}{\hyperlink{call-esq_119_018-3-1}{p.\,16}}
\DeclareNoteReturns{copy1-4}{\hyperlink{call-esq_119_019-4-1}{p.\,17}}
\DeclareNoteReturns{copy1-5}{\hyperlink{call-esq_119_026-5-1}{p.\,24}}
\DeclareNoteReturns{copy1-6}{\hyperlink{call-esq_119_033-6-1}{p.\,31}}
\DeclareNoteReturns{copy1-7}{\hyperlink{call-esq_119_036-7-1}{p.\,34}}
\begin{document}
\begin{titlepage}
\centering
\vspace*{18mm}
{\Large Alexander Grothendieck\par}
\vspace{12mm}
{\Huge\bfseries\EditionTitle\par}
\vspace{8mm}
{\Large\EditionForm\par}
\vspace{14mm}
{\large January 1984\par}
\vfill
{\large Université de Montpellier\par fonds Grothendieck, cote 119\par}
\vspace{8mm}
{\normalsize Two typescript copies: sites 2--59 and 78--135\par}
\vspace{8mm}
{\normalsize Edition dated 14 November 2026\par}
\end{titlepage}
\thispagestyle{plain}
\section*{Preface}
\noindent\textbf{Partial working edition.} This file contains 58 of the 58 source pages required for this form. Collation and final validation remain in progress.\par\smallskip
This English translation follows the corrected French text of copy 1. Authorial changes and justified corrections of isolated typing slips are recorded in the apparatus. Copy 2 annotations are identified separately. Translator’s notes explain words, wordplay, abbreviations and notation at the place where an explanation is needed. Unresolved readings remain marked and linked to their scans. The author’s notes (1)–(7) remain at the end, with links in both directions. The formulas retain the exact French source strings; terminology follows the supplied unified English glossary.

\smallskip
AI-transcribed by Claude Opus 5.5 (Anthropic, effort high, Message Batches API); checked against the scans, and this edition built, by GPT-6 Astra (OpenAI, ChatGPT Pro mode); prompts and records by Claude (Anthropic).
AI-translated by GPT-6 Astra (OpenAI, ChatGPT Pro mode).

The transcription of copy 1 was made with the text layer of the Grothendieck Circle's typeset \textit{Esquisse} (\texttt{EsquisseFr.pdf}) given to the transcribing model as a draft, to compare with after it had transcribed each page; every reading of this edition was decided on the scans.

The published French text and English translation in \textit{Geometric Galois Actions 1}, edited by L. Schneps and P. Lochak (Cambridge University Press, 1997), the Grothendieck Circle's typeset \texttt{EsquisseFr.pdf}, and the Centre for Grothendieckian Studies' draft transcription G94d are named as published texts. They are not reproduced here.

\section*{Rights and credits}
The rights in the original work belong to its author or the author's representatives or heirs. Where we hold rights, in the AI-produced transcription, translation, corrections and typesetting, they are dedicated to the public domain under CC0 1.0. This edition is published for educational, non-profit purposes, and we claim no rights we do not hold. If you hold rights in this work and object to its publication here, please open a GitHub issue and we will take it down.

\clearpage

% BEGIN PAGE esq_119_002
\scanpage{119/2}
\uline{SKETCH OF A PROGRAM}

by Alexandre Grothendieck

\uline{Contents} :

1. Opening.\tnote{TN-002-01}{The French heading is ``Envoi'', a word for a sending or dispatch, used here to introduce the text addressed to the application committee and other readers.}

2. A "Légo-Teichmüller game", and the Galois group of $\overline{\mathbb{Q}}$ over $\mathbb{Q}$.

3. Number fields associated with a child's drawing.\tnote{TN-002-02}{``Child's drawing'' translates ``dessin d'enfant''. The phrase becomes a mathematical term in the discussion.}

4. Regular polyhedra and finite fields.

5. A hue and cry against so-called ``general'' topology, and heuristic reflections
toward a so-called ``tame'' topology.\tnote{TN-002-03}{``Haro sur'' is a denunciatory cry directed against something; ``modérée'' is rendered ``tame'' throughout, following the terminology standard supplied for this edition.}

6. ``Differentiable theories'' (in the manner of Nash) and ``tame theories''.

7. Pursuing Stacks.\tnote{TN-002-04}{The French title is ``A la Poursuite des Champs''. ``Champ'' is translated ``stack'' in accordance with the supplied terminology standard.}

8. Digressions in two-dimensional geometry.

9. An appraisal of teaching activity.

10. Epilogue.

Notes.

N.B. The asterisks (*) refer to the notes appearing at the bottom of
the same page, the numbered references from (\textsuperscript{\hand{1}}) to (\textsuperscript{\hand{7}}\unsure{.}) to the notes (added
later) collected at the end of the report.

\mspage{\hand{2}}
\copytwo{A2-119-078-001}{\scancontext{https://grothendieck.umontpellier.fr/web/viewer.html?file=../119.pdf\#page=79}{\href{https://grothendieck.umontpellier.fr/web/viewer.html?file=../119.pdf\#page=79}{[copy 2, scan 119/78]} Copy 2 has in the upper-left margin: \unsure{45,49,53}. A dark mark obscures part of the digits.}}
\copytwo{A2-119-078-002}{\scancontext{https://grothendieck.umontpellier.fr/web/viewer.html?file=../119.pdf\#page=79}{\href{https://grothendieck.umontpellier.fr/web/viewer.html?file=../119.pdf\#page=79}{[copy 2, scan 119/78]} Copy 2 next has: \unsure{33,3\illegible{1}},26,8,. The illegibility mark here denotes an unresolved numerical token.}}
\copytwo{A2-119-078-003}{\scancontext{https://grothendieck.umontpellier.fr/web/viewer.html?file=../119.pdf\#page=79}{\href{https://grothendieck.umontpellier.fr/web/viewer.html?file=../119.pdf\#page=79}{[copy 2, scan 119/78]} Copy 2 has the handwritten numbers 1,2,9,10 on the next margin line.}}
\copytwo{A2-119-078-004}{\scancontext{https://grothendieck.umontpellier.fr/web/viewer.html?file=../119.pdf\#page=79}{\href{https://grothendieck.umontpellier.fr/web/viewer.html?file=../119.pdf\#page=79}{[copy 2, scan 119/78]} Copy 2 has the handwritten circled number \textcircled{17} in the margin.}}
\copytwo{A2-119-078-005}{\scancontext{https://grothendieck.umontpellier.fr/web/viewer.html?file=../119.pdf\#page=79}{\href{https://grothendieck.umontpellier.fr/web/viewer.html?file=../119.pdf\#page=79}{[copy 2, scan 119/78]} Copy 2 has handwritten archive foliation 78.}}

\finishscan
% END PAGE esq_119_002
% BEGIN PAGE esq_119_003
\scanpage{119/3}
\mspage{1}

SKETCH OF A PROGRAM

by A. Grothendieck

1. As the present circumstances make the prospects of teaching at research level
at the University increasingly illusory for me, I have resolved to apply for
admission to the CNRS,\tnote{TN-003-01}{CNRS: \href{https://aissai.cnrs.fr/en/legal-notice/}{Centre national de la recherche scientifique}, the French National Center for Scientific Research.}
so that I may devote my energy to developing work and perspectives
that, it is becoming clear, no student (nor even, it seems,
any fellow mathematician) will be found to develop in my
stead.

In place of a ``Qualifications and Works'' document,\tnote{TN-003-02}{``Titres et Travaux'' is the French designation of the document listing qualifications and scholarly work.}
following this text will be found the full reproduction of a sketch, arranged by themes,
of what I regarded as my main mathematical contributions at the time
of writing that report, in 1972. It also contains
a list of articles published by that date. I have published no scientific
articles since 1970. In the lines that follow,
I propose to give at least a glimpse of some principal themes
of my mathematical reflections since then. Over the years these reflections
have taken material form in two bulky boxes of
manuscript notes, probably difficult for anyone other than
myself to decipher, which, after successive stages of decantation,
are perhaps awaiting their moment for an overall write-up, at least
a provisional one, intended for the mathematical community. The
term ``write-up'' here is somewhat inappropriate,\tnote{TN-003-03}{``Write-up'' renders ``rédaction''. The source itself immediately explains the continuing work of discovery meant by this word.}
since it is much more a matter of developing the many ideas and visions begun
over these last twelve years, making them more precise and exploring them
more deeply, with all the unforeseen turns that constantly
accompany this kind of work - a work of discovery, then, and
not of compiling notes piously accumulated. And I fully intend,
in writing the ``Mathematical Reflections'' begun in February
1983,\tnote{TN-003-04}{The French title is ``Réflexions Mathématiques''.}
to let the course of thought that probes and discovers appear
clearly through the pages, often groping in the half-light,
with sudden openings of light when some stubborn

\margin{3}
\copytwo{A2-119-079-001}{\scancontext{https://grothendieck.umontpellier.fr/web/viewer.html?file=../119.pdf\#page=80}{\href{https://grothendieck.umontpellier.fr/web/viewer.html?file=../119.pdf\#page=80}{[copy 2, scan 119/79]} Copy 2 has the title ‘ESQUISSE D'UN PROGRAMME’ crossed out with oblique strokes.}}
\copytwo{A2-119-079-002}{\scancontext{https://grothendieck.umontpellier.fr/web/viewer.html?file=../119.pdf\#page=80}{\href{https://grothendieck.umontpellier.fr/web/viewer.html?file=../119.pdf\#page=80}{[copy 2, scan 119/79]} Copy 2 has the line ‘par A. Grothendieck’ crossed out with the same oblique strokes.}}
\copytwo{A2-119-079-003}{\scancontext{https://grothendieck.umontpellier.fr/web/viewer.html?file=../119.pdf\#page=80}{\href{https://grothendieck.umontpellier.fr/web/viewer.html?file=../119.pdf\#page=80}{[copy 2, scan 119/79]} Copy 2 has the handwritten foliation 79 at the lower left.}}

\finishscan
% END PAGE esq_119_003
% BEGIN PAGE esq_119_004
\scanpage{119/4}
\mspage{2}

false, or merely inadequate\frapparatus{V-FR-004-001}, image is finally flushed out
and brought to light, and the things that seemed askew
fall into place, in the mutual harmony proper to them.
\qquad Be that as it may, the following sketch of some
themes of reflection from the last ten or twelve years will
serve at the same time as a sketch of a work program for the
coming years, which I intend to devote to developing
these themes, or at least some of them. It is intended,
on the one hand, for the colleagues on the National Committee called upon to rule
on my application,\tnote{TN-004-01}{``National Committee'' translates ``Comité National'', the committee named in the source in connection with the CNRS application.}
and, on the other, for a few other colleagues, former
students, and friends, in the event that some of the ideas sketched
here might interest one of them.

2. \qquad The demands of university teaching, addressed
therefore to students (including so-called ``advanced'' students) whose
mathematical background is modest (and often less than modest), have led me
to renew drastically the themes of reflection to propose
to my students, and, one thing leading to another and increasingly, to myself
as well. It had seemed important to me to start from a common
intuitive background, independent of any technical language meant to express it,
far earlier than any such language - it turned out that the geometric
and topological intuition of shapes, and more particularly of
two-dimensional shapes, was such a common ground. Thus the themes in question
can be grouped under the name ``topology of surfaces'' or
``geometry of surfaces'', with the understanding, in the latter name,
that the principal emphasis falls on the topological properties
of surfaces, or on the combinatorial aspects that constitute
their most down-to-earth technical expression, and not on the differential
aspects, or even the conformal, riemannian, holomorphic and (thence)
the ``complex algebraic curves'' aspect. Once this last step
has been taken, however, suddenly algebraic geometry (my old
love !) bursts in again, through the very objects that
can be regarded as the ultimate building blocks of all
the other algebraic varieties. Whereas in my research before

\margin{4}
\copytwo{A2-119-080-002}{\scancontext{https://grothendieck.umontpellier.fr/web/viewer.html?file=../119.pdf\#page=81}{\href{https://grothendieck.umontpellier.fr/web/viewer.html?file=../119.pdf\#page=81}{[copy 2, scan 119/80]} Copy 2 has the handwritten foliation 80 at the lower left.}}

\finishscan
% END PAGE esq_119_004
% BEGIN PAGE esq_119_005
\scanpage{119/5}
\mspage{3}

1970, my attention was systematically directed toward objects
of the greatest generality, in order to draw out an adequate overall language
for the world of algebraic geometry, and I lingered
over algebraic curves only to the strict extent that this proved
indispensable (particularly in étale cohomology) for developing
``all-purpose'' techniques and statements valid in every dimension
and in every place (I mean, over all base schemes, or even
all ringed base topoi ...), here I am thus brought back, through
objects so simple that a child can come to know them through play,
to the beginnings and origins of algebraic geometry, familiar to Riemann
and his followers !

Since about 1975, then, the geometry of (real)
surfaces, and from 1977 onward the connections between questions in
the geometry of surfaces and the algebraic geometry of algebraic curves
defined over fields such as $\mathbb{C}$, $\mathbb{R}$ or extensions of finite
type of $\mathbb{Q}$, have been my principal source of inspiration, as well
as my constant guiding thread. It was with surprise and wonder
that over the years I discovered (or rather, probably,
rediscovered) the prodigious, truly inexhaustible richness, the
unsuspected depth of this seemingly so innocuous theme. I believe
I sense in it a particularly vital point, a privileged point of convergence
of the main currents of mathematical ideas, as also
of the main structures and the visions of things they express,
from the most specific (such as the rings $\mathbb{Z}$, $\mathbb{Q}$, $\bar{\mathbb{Q}}$, $\mathbb{R}$, $\mathbb{C}$,
or the group $S\ell(2)$ over one of these rings, or general reductive
algebraic groups) to the most ``abstract'', such as algebraic,
complex analytic or real analytic ``multiplicities''.\tnote{TN-005-01}{``Multiplicity'' consistently renders the author's ``multiplicité''. His parenthesis here explains its role in retaining the automorphisms of the objects under study.}
(These enter naturally when one sets out to study
``moduli varieties'' systematically for the geometric objects
under consideration, seeking to go beyond the notoriously insufficient
viewpoint of ``coarse moduli'', which amounts, most unfortunately,
to ``killing'' the automorphism groups of these objects .) Among these
moduli multiplicities, it is those of Mumford-Deligne for
``stable'' algebraic curves of genus $g$, with $\nu$ marked points, which

\margin{5}
\copytwo{A2-119-081-001}{\scancontext{https://grothendieck.umontpellier.fr/web/viewer.html?file=../119.pdf\#page=82}{\href{https://grothendieck.umontpellier.fr/web/viewer.html?file=../119.pdf\#page=82}{[copy 2, scan 119/81]} Copy 2 has the handwritten foliation 81 at the lower left.}}

\finishscan
% END PAGE esq_119_005
% BEGIN PAGE esq_119_006
\scanpage{119/6}
\mspage{4}

I denote by $M^{\wedge}_{g,\nu}$ (the compactification of the ``open'' multiplicity
$M_{g,\nu}$ corresponding to smooth curves), that for some two
or three years have exerted a particular fascination upon me,
perhaps stronger than that of any other mathematical object to date.
To tell the truth, it is rather the system of \uline{all} the multiplicities
$M_{g,\nu}$ for varying $g,\nu$, connected with one another by a certain number
of fundamental operations (such as the operations of ``filling
holes'', i.e. ``erasing'' marked points, that of ``gluing'',
and the inverse operations), which reflect, in absolute algebraic
geometry of characteristic zero (for the moment), familiar geometric
operations from the viewpoint of topological or conformal ``surgery''
on surfaces.\frapparatus{V-FR-006-001} The principal reason for this fascination is probably
that this very rich geometric structure on the
system of ``open'' moduli multiplicities $M_{g,\nu}$ is reflected
in an analogous structure on the corresponding fundamental
groupoids, the ``Teichmüller groupoids'' $\hat{T}_{g,\nu}$ ,and that these operations
at the level of the $\hat{T}_{g,\nu}$ have a sufficiently intrinsic character
for the Galois group $\mathrm{I}\!\Gamma$ of $\overline{\mathbb{Q}}/\mathbb{Q}$ to act on this entire ``tower'' of
Teichmüller groupoids, respecting all these structures. More
extraordinary still, this action is \uline{faithful} - in fact,
it is already faithful on the first nontrivial ``level'' of this tower,
namely $\hat{T}_{0,4}$ - which also means, essentially, that the outer
action of $\mathrm{I}\!\Gamma$ on the fundamental group $\hat{\pi}_{0,3}$ of the standard
projective line $\mathbb{P}^1$ over $\mathbb{Q}$, with the three points $0,\ 1,\ \infty$ removed, is
already faithful. Thus \uline{the Galois group} $\mathrm{I}\!\Gamma$ \uline{is realized as a group}
\uline{of automorphisms of a most concrete profinite group}\frapparatus{V-FR-006-002}, moreover respecting
certain essential structures of that group. It follows
that an element of $\mathrm{I}\!\Gamma$ can be ``parametrized'' (in various equivalent
ways, moreover) by a suitable element of this profinite group
$\hat{\pi}_{0,3}$ (a free profinite group on two generators), or by a system
of such elements, this element or these elements being, moreover, subject
to certain simple conditions, necessary (and probably insufficient)
for this element or these elements actually to correspond to an element of $\mathrm{I}\!\Gamma$.
One of the most fascinating tasks here is precisely to grasp
necessary \uline{and} sufficient conditions on an outer automorphism

\margin{6}
\copytwo{A2-119-082-001}{\scancontext{https://grothendieck.umontpellier.fr/web/viewer.html?file=../119.pdf\#page=83}{\href{https://grothendieck.umontpellier.fr/web/viewer.html?file=../119.pdf\#page=83}{[copy 2, scan 119/82]} Copy 2 has the handwritten foliation 82 at the lower left.}}

\finishscan
% END PAGE esq_119_006
% BEGIN PAGE esq_119_007
\scanpage{119/7}
\mspage{5}

of $\hat{\pi}_{0,3}$, i.e. on the corresponding parameter or parameters, for it
to come from an element of $\mathrm{I}\!\Gamma$ - which would provide a
``purely algebraic'' description, in terms of profinite groups and without
reference to the Galois theory of number fields, of the Galois
group $\mathrm{I}\!\Gamma = \mathrm{Gal}(\bar{\mathbb{Q}}/\mathbb{Q})$ !

Perhaps even a conjectural characterization of $\mathrm{I}\!\Gamma$
as a subgroup of Autext($\hat{\pi}_{0,3}$) is for the moment out of
reach\notecall{copy1-1}{call-esq_119_007-1-1}{($^1$)} ; I have no conjecture to propose yet. Another
task, however, can be approached immediately: describing
the action of $\mathrm{I}\!\Gamma$ on the entire Teichmüller tower in terms of its
action on the ``first level'' $\hat{\pi}_{0,3}$ , i.e. expressing an automorphism
of this tower in terms of the ``parameter'' in $\hat{\pi}_{0,3}$ that specifies
the varying element $\gamma$ of $\mathrm{I}\!\Gamma$. This is connected with a representation of the
Teichmüller tower (as a groupoid equipped with a ``gluing''
operation) by generators and relations, which will in particular give
a presentation by generators and relations, in the ordinary sense, of
each of the $\hat{T}_{g,\nu}$ (as a profinite groupoid). Here, even for
$g = 0$ (thus when the corresponding Teichmüller groups are
``well-known'' braid groups), the generators and relations
known to date that have come to my attention seem unusable
as they stand, because they lack the invariance and
symmetry properties indispensable for the action of $\mathrm{I}\!\Gamma$ to be directly
readable from this presentation. This is connected in particular with the fact that
people still persist, when calculating with fundamental
groups, in fixing a single base point, instead of cleverly choosing
a whole collection invariant under the symmetries of the
situation, which are thereby lost along the way. In certain situations
(such as descent theorems in the manner of Van Kampen for fundamental
groups), it is much more elegant, even indispensable\frapparatus{V-FR-007-001} for understanding
anything about them, to work with fundamental groupoids
relative to a suitable collection of base points , and this is certainly
so for the Teichmüller tower. It would seem (incredible,
but true !) that the very geometry of the first level of the
Teichmüller tower (thus corresponding to ``moduli'' either for projective

\margin{7}
\copytwo{A2-119-083-001}{\scancontext{https://grothendieck.umontpellier.fr/web/viewer.html?file=../119.pdf\#page=84}{\href{https://grothendieck.umontpellier.fr/web/viewer.html?file=../119.pdf\#page=84}{[copy 2, scan 119/83]} Copy 2 has the handwritten foliation 83 at the lower left.}}

\finishscan
% END PAGE esq_119_007
% BEGIN PAGE esq_119_008
\scanpage{119/8}
\mspage{6}

lines with four marked points or for elliptic
curves (!)) has never been properly made explicit, for example the relation
between the genus 0 case and the geometry of the octahedron and of the
tetrahedron. A fortiori, the moduli multiplicities $M_{0,5}$ (for
projective lines with five marked points) and $M_{1,2}$ (for
curves of genus 1 with two marked points), which are moreover almost
isomorphic to one another, seem virgin territory - braid groups
will not enlighten us about them ! I have begun to
look at $M_{0,5}$ in spare moments; it is a true jewel, with a
very rich geometry closely connected with that of the icosahedron.

The a priori interest of a complete knowledge of the first two
levels of the tower (namely, the cases in which the moduli dimension
$N = 3g-3+\nu$ is $\leq 2$) lies in this principle, that \uline{the entire tower
is reconstructed from the first two levels}, in the sense that, through
the fundamental ``gluing'' operation, level 1 supplies a
complete system of generators, and level 2 a complete system of
relations. There is a striking analogy, and I am convinced it is
more than formal, between this principle and Demazure's analogous
principle for the structure of reductive algebraic groups, when
the term ``level'' or ``moduli dimension'' is replaced by ``semisimple
rank of the reductive group''. The connection becomes still more striking
when one recalls that the Teichmüller group $T_{1,1}$ (now in the
discrete transcendental context, and not in the profinite algebraic
context, where one finds the profinite completions of the former) is
precisely $S\ell(2,\mathbb{Z})$, i.e. the group of integral points of the
``absolute'' simple group scheme of rank 1 $S\ell(2)_{\mathbb{Z}}$. Thus, \uline{the fundamental
building block for the Teichmüller tower} \uline{is essentially
the same as that for ``the tower'' of reductive groups of all ranks} -
a group of which one can probably say, moreover, that it is present
in all the essential disciplines of mathematics.

This principle for constructing the Teichmüller tower has
not been proved at present - but I have no doubt that it
is valid. It would follow (through a theory of dévissage of stratified
structures - in this instance the $M^{\wedge}_{g,\nu}$ - still to be written,
cf. par. 5) from an extremely plausible property of the open

\margin{8}
\copytwo{A2-119-084-001}{\scancontext{https://grothendieck.umontpellier.fr/web/viewer.html?file=../119.pdf\#page=85}{\href{https://grothendieck.umontpellier.fr/web/viewer.html?file=../119.pdf\#page=85}{[copy 2, scan 119/84]} Copy 2 has the handwritten foliation 84 at the lower left.}}

\finishscan
% END PAGE esq_119_008
% BEGIN PAGE esq_119_009
\scanpage{119/9}
\mspage{7}

\unsure{moduli}\tnote{TN-009-01}{The checked French retains the alternatives ``mcdulaires|modulaires'' for the unclear typed letter. Both lead to ``moduli'' in this phrase; the doubt remains marked and linked to the scan.}
multiplicities $M_{g,\nu}$ in the complex analytic context,
namely that, for a moduli dimension $N \geq 3$, the fundamental
group of $M_{g,\nu}$ (i.e. the usual Teichmüller group $T_{g,\nu}$)
is isomorphic to the ``fundamental group at infinity'', i.e. that of a
``tubular neighborhood of infinity''. This is a very familiar fact
(due essentially to Lefschetz) for a smooth \uline{affine} variety of
dimension $N \geq 3$. It is true that the moduli multiplicities are
not affine (except for small values of $g$), but it would suffice
for such an $M_{g,\nu}$ of dimension $N$ (or rather, a suitable finite
covering)\frapparatus{V-FR-009-001} to be the union of $N$-2 affine open sets, thus
for $M_{g,\nu}$ not to be ``too close to a compact variety''.

Having no doubt about this principle for constructing the
Teichmüller tower, I prefer to leave it to the experts in the transcendental
theory, better equipped than I am, to prove what is needed
(should any be interested), and instead to make explicit,
with all the care it deserves, the resulting structure of the
Teichmüller tower by generators and relations, this time in the
discrete setting and not the profinite one - which amounts, essentially,
to a complete understanding of the four moduli multiplicities
$M_{0,4}$, $M_{1,1}$, $M_{0,5}$, $M_{1,2}$ , and of their fundamental groupoids relative
to suitably chosen ``base points''. These present themselves
quite naturally as the complex algebraic curves of the type
$(g,\nu)$ under consideration whose automorphism group (necessarily
finite) is larger than in the generic case (*). Including among them the
holomorphic sphere with three marked points (coming from $M_{0,3}$, i.e.

\rule{4cm}{0.4pt}

(*) One must also add the ``base points'' arising from operations
of gluing ``pieces'' of the same type in lower moduli dimension.
On the other hand, in moduli dimension 2 (the cases of $M_{0,5}$ and $M_{1,2}$),
one should exclude the points of certain one-parameter families
of curves admitting an exceptional automorphism of order 2. These
families, moreover, constitute remarkable rational curves on the multiplicities
under consideration, which seem to me an important ingredient
of the structure of these multiplicities.

\margin{9}
\copytwo{A2-119-085-002}{\scancontext{https://grothendieck.umontpellier.fr/web/viewer.html?file=../119.pdf\#page=86}{\href{https://grothendieck.umontpellier.fr/web/viewer.html?file=../119.pdf\#page=86}{[copy 2, scan 119/85]} Copy 2 has a blue proof sign in the left margin, approximately \unsure{$\frown$}, beside ‘convena\deleted{n}ble’.}}
\copytwo{A2-119-085-003}{\scancontext{https://grothendieck.umontpellier.fr/web/viewer.html?file=../119.pdf\#page=86}{\href{https://grothendieck.umontpellier.fr/web/viewer.html?file=../119.pdf\#page=86}{[copy 2, scan 119/85]} Copy 2 has the handwritten foliation 85 at the lower left.}}

\finishscan
% END PAGE esq_119_009
% BEGIN PAGE esq_119_010
\scanpage{119/10}
\mspage{8}

from level O), one finds \uline{twelve} \uline{fundamental} ``\uline{building blocks}''
(6 of genus O, 6 of genus 1) in a ``Légo-Teichmüller game (large
box), where the marked points on the surfaces under consideration are replaced
by ``holes'' with boundary, so as to obtain surfaces with boundary,
thus building blocks that can be fitted together with gentle
friction, as in the ordinary Légo game dear to our children
(or grandchildren...). Assembling them provides a thoroughly
visual way of constructing every type of surface (it is essentially
these assemblies that will be the ``base points'' for
our famous tower), and also of visualizing the \uline{elementary} ``\uline{paths}''
through equally concrete operations, such as ``twists'', or
automorphisms of the pieces in the game, and of writing the \uline{fundamental relations}
between composite paths. Depending on the size (and the price !) of the
construction box used, one moreover \unsure{finds}\tnote{TN-010-01}{The French reading is retained as ``trcuve|trouve'' because the letterform is unresolved. Both alternatives are represented here by ``finds'', with the source doubt retained.}
many different descriptions of the Teichmüller tower by generators and
relations. The smallest box is reduced to pieces that are all
identical, of type (O,3) - these are Thurston's ``pairs of pants'',\tnote{TN-010-02}{``Pairs of pants'' translates ``pantalons'', the surface name used by the author. The letter O in the scanned type expressions is retained.}
and the Légo-Teichmüller game I am trying to describe, arising from
motivations and reflections in absolute algebraic geometry over the field $\mathbb{Q}$,
is very close to Thurston's game of ``hyperbolic geodesic surgery'',
whose existence I learned of last year from Yves Ladegaillerie.
In a microseminar with Carlos Contou-Carrère and Yves \var{Ladegraillerie}{Ladegaillerie}{slip corrected}{V-FR-010-001},
we have begun a reflection one of whose aims is to confront
the two viewpoints, which complement each other.

I add that each of the twelve building blocks in
the ``large box'' is equipped with a canonical cellular decomposition\frapparatus{V-FR-010-002},
stable under all symmetries, whose only vertices are the
``marked points'' (or centers of the holes), and whose edges are certain
geodesic paths (for the canonical riemannian structure on the sphere
or torus under consideration) between certain pairs of vertices (namely those
that lie on the same ``real locus'', for a suitable real structure
on the complex algebraic curve under consideration). Consequently,
all the surfaces obtained in this game by assembly are equipped
with canonical cellular structures, which in their turn (cf.§3 below)

\margin{10}
\copytwo{A2-119-086-001}{\scancontext{https://grothendieck.umontpellier.fr/web/viewer.html?file=../119.pdf\#page=87}{\href{https://grothendieck.umontpellier.fr/web/viewer.html?file=../119.pdf\#page=87}{[copy 2, scan 119/86]} Copy 2 has a handwritten sign \unsure{*} after ‘Teichmüller’; its exact form and any possible covered letter or punctuation remain uncertain.}}
\copytwo{A2-119-086-002}{\scancontext{https://grothendieck.umontpellier.fr/web/viewer.html?file=../119.pdf\#page=87}{\href{https://grothendieck.umontpellier.fr/web/viewer.html?file=../119.pdf\#page=87}{[copy 2, scan 119/86]} Copy 2 has a handwritten sign \unsure{*} in the right margin on the same line.}}
\copytwo{A2-119-086-003}{\scancontext{https://grothendieck.umontpellier.fr/web/viewer.html?file=../119.pdf\#page=87}{\href{https://grothendieck.umontpellier.fr/web/viewer.html?file=../119.pdf\#page=87}{[copy 2, scan 119/86]} Copy 2 circles the typed words ‘dans ce jeu’ and connects them by a line to the position after ‘Par suite,’; the words are not crossed out.}}
\copytwo{A2-119-086-004}{\scancontext{https://grothendieck.umontpellier.fr/web/viewer.html?file=../119.pdf\#page=87}{\href{https://grothendieck.umontpellier.fr/web/viewer.html?file=../119.pdf\#page=87}{[copy 2, scan 119/86]} Copy 2 has the handwritten foliation 86 at the lower left.}}

\finishscan
% END PAGE esq_119_010
% BEGIN PAGE esq_119_011
\scanpage{119/11}
\mspage{9}

allow these surfaces to be considered as associated with
canonically determined complex algebraic curves (and even curves over $\overline{\mathbb{Q}}$).
Here there is a typical interplay between the combinatorial
and the complex algebraic (or better, the algebraic over $\overline{\mathbb{Q}}$).

The ``small box'', with its identical pieces, which has
the charm of economy, will probably give a relatively
complicated description of the relations (complicated, but by no means
inextricable !). The large box will give rise to more numerous
relations (because there are many more base points and
remarkable paths between them), but with a more transparent structure. I
expect that in moduli dimension 2, just as in the more or
less familiar case of moduli dimension 1 (with, in particular, the
description of $S\ell(2,\mathbb{Z})$ by $(\rho,\sigma \mid \rho^3 = \sigma^2,\ \rho^4 = \sigma^6 = 1)$), one will find
generation by the automorphism groups of the three relevant types
of pieces, with simple relations that I have not worked out
at the time of writing these lines. Perhaps one will even find a
principle of this kind for all the $T_{g,\nu}$, as well as a cellular
decomposition of $M^{\wedge}_{g,\nu}$ generalizing those that present themselves spontaneously
for $M^{\wedge}_{0,4}$ and $M^{\wedge}_{1,1}$, and that I already glimpse for
moduli dimension 2, using the hypersurfaces corresponding
to the various \uline{real structures} on the complex structures
under consideration to carry out the desired cellular decomposition.


3. Instead of following (as planned) a rigorous thematic
order, I have let myself be carried away by my predilection for a
particularly rich and burning theme, to which I moreover intend to
devote myself as a priority for some time, beginning with the start
of the 84/85 academic year. I therefore resume the thematic account where I left it,
at the very beginning of the preceding paragraph.

My interest in topological surfaces begins to emerge
in 1974, when I propose to Yves Ladegaillerie the theme of studying
isotopically the embeddings of a topological 1-complex in a compact
surface. In the two years that follow, this study leads him to
a remarkable isotopy theorem, giving a complete algebraic description

\margin{11}
\copytwo{A2-119-087-001}{\scancontext{https://grothendieck.umontpellier.fr/web/viewer.html?file=../119.pdf\#page=88}{\href{https://grothendieck.umontpellier.fr/web/viewer.html?file=../119.pdf\#page=88}{[copy 2, scan 119/87]} Copy 2 has $\over{\rho}{\sigma}^4 = \over{\sigma}{\rho}^6 = 1$: handwritten $\sigma$ and $\rho$ cover the typed letters $\rho$ and $\sigma$ respectively.}}
\copytwo{A2-119-087-002}{\scancontext{https://grothendieck.umontpellier.fr/web/viewer.html?file=../119.pdf\#page=88}{\href{https://grothendieck.umontpellier.fr/web/viewer.html?file=../119.pdf\#page=88}{[copy 2, scan 119/87]} Copy 2 has two oblique strokes // in the left margin beside the overwritten formula.}}
\copytwo{A2-119-087-003}{\scancontext{https://grothendieck.umontpellier.fr/web/viewer.html?file=../119.pdf\#page=88}{\href{https://grothendieck.umontpellier.fr/web/viewer.html?file=../119.pdf\#page=88}{[copy 2, scan 119/87]} Copy 2 has a small point above the e in ‘Yv\unsure{ė|e}s’; its origin remains uncertain.}}
\copytwo{A2-119-087-004}{\scancontext{https://grothendieck.umontpellier.fr/web/viewer.html?file=../119.pdf\#page=88}{\href{https://grothendieck.umontpellier.fr/web/viewer.html?file=../119.pdf\#page=88}{[copy 2, scan 119/87]} Copy 2 has the handwritten foliation 87 at the lower left.}}

\finishscan
% END PAGE esq_119_011
% BEGIN PAGE esq_119_012
\scanpage{119/12}
\mspage{10}

of the isotopy classes of embeddings of such 1-complexes,
or of compact surfaces with boundary, in a compact oriented surface,
in terms of certain very simple combinatorial invariants and the
fundamental groups of the participants. This theorem, which should
extend without difficulty to embeddings of any compact space
(triangulable, to simplify matters) in a compact oriented surface,
recovers as easy corollaries several deep classical results
of surface theory, and in particular Baer's isotopy theorem.
It is finally going to be published, separately from the rest (and
ten years later, given the harshness of the times...), in Topology. Ladegaillerie's
work also contains a purely algebraic description,
in terms of fundamental groupoids, of the
``isotopy'' category of compact surfaces $X$, equipped with a topological
1-complex $K$ embedded in $X$. This description, which had the misfortune
to go against the ``taste of the day'' and consequently seems unpublishable,
has nevertheless served (and still serves) as a valuable guide in
my subsequent reflections, particularly in the context of absolute
algebraic geometry of characteristic zero.

The case in which $(X,K)$ is a 2-dimensional ``map'',\tnote{TN-012-01}{``Map'' here translates ``carte'', a surface equipped with the embedded graph and cell conditions specified in the sentence.}
i.e. in which the connected components of $X\backslash K$ are open 2-cells \unsure{.}
(and in which $K$ is moreover equipped with a finite set $S$ of ``vertices'',
such that the connected components of $K\backslash S$ are open 1-cells)
gradually attracts my attention in the following years.
The isotopy category of these maps admits a particularly
simple algebraic description, through the set of ``frames'' (or ``flags'',
or ``biarcs'') associated with the map, which is naturally equipped
with the structure of a set with a group of operators, under the group
\[ \underline{C}_2 = \langle \sigma_0,\sigma_1,\sigma_2 \,|\, \sigma_0^{\,2} = \sigma_1^{\,2} = \sigma_2^{\,2} = (\sigma_0\sigma_2)^2 = 1 \rangle \ , \]
which I call the (unoriented) \uline{cartographic group} of dimension 2.
It admits\frapparatus{V-FR-012-001} as a subgroup of index 2 the \uline{oriented cartographic group},
generated by products of even numbers of the generators, which can
also be described as
\[ \underline{C}_2^{+} = \langle \rho_s,\rho_f,\sigma \,|\, \rho_s\rho_f = \sigma,\ \sigma^2 = 1 \rangle \ , \]

\margin{12}
\copytwo{A2-119-088-001}{\scancontext{https://grothendieck.umontpellier.fr/web/viewer.html?file=../119.pdf\#page=89}{\href{https://grothendieck.umontpellier.fr/web/viewer.html?file=../119.pdf\#page=89}{[copy 2, scan 119/88]} Copy 2 has ‘i\over{m}{n}publia-ble’: a handwritten n covers the typed m.}}
\copytwo{A2-119-088-002}{\scancontext{https://grothendieck.umontpellier.fr/web/viewer.html?file=../119.pdf\#page=89}{\href{https://grothendieck.umontpellier.fr/web/viewer.html?file=../119.pdf\#page=89}{[copy 2, scan 119/88]} Copy 2 also has a handwritten n in the right margin of this line.}}
\copytwo{A2-119-088-003}{\scancontext{https://grothendieck.umontpellier.fr/web/viewer.html?file=../119.pdf\#page=89}{\href{https://grothendieck.umontpellier.fr/web/viewer.html?file=../119.pdf\#page=89}{[copy 2, scan 119/88]} Copy 2 has the handwritten foliation 88 at the lower left, preceded by a point of uncertain origin: \unsure{.}88.}}

\finishscan
% END PAGE esq_119_012
% BEGIN PAGE esq_119_013
\scanpage{119/13}
\mspage{11}

(with
\[ \rho_s = \sigma_2\sigma_1 \ , \quad \rho_f = \sigma_1\sigma_0 \ , \quad \sigma = \sigma_0\sigma_2 = \sigma_2\sigma_0 \ , \]
the operations of \uline{elementary rotation} of a frame about a vertex,
a face and an edge respectively). There is a perfect dictionary
between the topological situation of compact maps, resp.
compact oriented maps, on the one hand, and finite sets with
operator group $\underline{C}_2$ resp. $\underline{C}_2^+$ on the other, a dictionary whose
existence was, moreover, more or less known, but never
stated with the necessary precision or developed to any extent.
This foundational work is carried out with the care it deserves in
an excellent DEA project,\tnote{TN-013-01}{DEA abbreviates ``diplôme d'études approfondies'', an advanced studies diploma.}
undertaken jointly by Jean Malgoire and
Christine Voisin in 1976.

This reflection suddenly takes on a new dimension
with the simple observation that the group $\underline{C}_2^+$ can be interpreted
as a quotient of the fundamental group of an oriented sphere with
three points removed, numbered 0, 1, 2, the operations $\rho_s$, $\sigma$, $\rho_f$
being interpreted as the loops about these points, satisfying
the familiar relation
\[ \ell_0\ell_1\ell_2 = 1 \ , \]
while the additional relation $\sigma^2 = 1$, i.e. $\ell_1^{\,2} = 1$, means
that one is interested in the quotient of the fundamental group corresponding
to an imposed ramification index 2 at the point 1, which therefore classifies
coverings of the sphere ramified at most at the
points 0, 1, 2, with ramification equal to 1 or 2 at the points
above 1. Thus compact oriented maps form an
isotopy category equivalent to that of these coverings, subject
to the further condition of being finite coverings.
Taking now as the reference sphere the Riemann sphere,
or complex projective line, rigidified by the three points 0, 1\frapparatus{V-FR-013-001}
and $\infty$ (the last thus replacing 2), and recalling that every
finite ramified covering of a complex algebraic curve itself inherits
the structure of a complex algebraic curve, one arrives at this
observation, which eight years later still seems to me just as
extraordinary: \uline{every oriented} ``\uline{finite}'' \uline{map is canonically realized}
\uline{on a complex algebraic curve} ! Better still, since the complex projective

\margin{13}
\copytwo{A2-119-089-001}{\scancontext{https://grothendieck.umontpellier.fr/web/viewer.html?file=../119.pdf\#page=90}{\href{https://grothendieck.umontpellier.fr/web/viewer.html?file=../119.pdf\#page=90}{[copy 2, scan 119/89]} Copy 2 has the handwritten foliation 89 at the lower left.}}

\finishscan
% END PAGE esq_119_013
% BEGIN PAGE esq_119_014
\scanpage{119/14}
\mspage{12}

line is defined over the absolute base field $\mathbb{Q}$ ,
as are the allowed ramification points, the algebraic curves
obtained are defined not only over $\mathbb{C}$, but over the algebraic
closure $\overline{\mathbb{Q}}$ of $\mathbb{Q}$ in $\mathbb{C}$. As for the original map, it
is recovered on the algebraic curve as the inverse image of the real
segment $[0,1]$ (where $0$ is regarded as a vertex, and $1$ as
the midpoint of a ``folded edge'' with 1 as its center), which constitutes
the ``universal oriented 2-map'' on the Riemann sphere (*).
The points of the algebraic curve $X$ above 0, 1 and $\infty$
are precisely the vertices and the ``centers'' of the edges and
faces respectively of the map $(X,K)$, and the orders of the vertices
and faces are precisely the multiplicities of the zeros and
poles of the rational function (defined over $\overline{\mathbb{Q}}$) on $X$,
expressing its structural projection to $\mathbb{P}^1_{\mathbb{C}}$.

This discovery, which technically amounts to so little,
made a very strong impression on me, and represents
a decisive turning point in the course of my reflections, in particular
a shift in my center of interest in mathematics, which
suddenly became sharply localized. I do not believe that a
mathematical fact has ever struck me as much as this one, or had
a comparable psychological impact \notecall{copy1-2}{call-esq_119_014-2-1}{($^2$)}. This surely has to do with the
so familiar, nontechnical nature of the objects under consideration,
of which any child's drawing scribbled on a scrap of paper (provided
the drawing is all in one piece) gives a perfectly
explicit example. With such a drawing are associated subtle arithmetic
invariants, which will be completely upset as soon as
one adds another stroke. Since we are dealing here with spherical maps, necessarily

\noindent\rule{4cm}{0.4pt}

(*) There is an analogous description of finite unoriented maps,
possibly with boundary, in terms of \uline{real} algebraic curves,
more precisely of a covering of $\mathbb{P}^1_{\mathbb{R}}$ ramified only at 0, 1, $\infty$\frapparatus{V-FR-014-001},
the surface with boundary associated with such a covering being $X(\mathbb{C})/\tau$, where $\tau$
is complex conjugation. The ``universal'' unoriented map is
here the disk, or upper hemisphere of the Riemann sphere, equipped,
as before, with the embedded 1-complex $K = [0,1]$.

\margin{14}
\copytwo{A2-119-090-001}{\scancontext{https://grothendieck.umontpellier.fr/web/viewer.html?file=../119.pdf\#page=91}{\href{https://grothendieck.umontpellier.fr/web/viewer.html?file=../119.pdf\#page=91}{[copy 2, scan 119/90]} Copy 2 has \inserted{re} above the beginning of “trouvé”, producing “retrouvé”.}}
\copytwo{A2-119-090-002}{\scancontext{https://grothendieck.umontpellier.fr/web/viewer.html?file=../119.pdf\#page=91}{\href{https://grothendieck.umontpellier.fr/web/viewer.html?file=../119.pdf\#page=91}{[copy 2, scan 119/90]} Copy 2 has two small points immediately to the right of the comma after “disque”; their origin and textual value remain unresolved.}}
\copytwo{A2-119-090-003}{\scancontext{https://grothendieck.umontpellier.fr/web/viewer.html?file=../119.pdf\#page=91}{\href{https://grothendieck.umontpellier.fr/web/viewer.html?file=../119.pdf\#page=91}{[copy 2, scan 119/90]} Copy 2 has the handwritten foliation 90 at the lower left.}}

\finishscan
% END PAGE esq_119_014
% BEGIN PAGE esq_119_015
\scanpage{119/15}
\mspage{13}

giving rise to curves of genus zero (which therefore
provide no ``moduli''), one can say that the curve
in question is ``pinned down'' as soon as three of its points are fixed, for
example three vertices of the map, or more generally three
centers of facets (vertices, edges or faces) - from that point on, the structural
map $f : X \to \mathbb{P}^{1}_{\mathbb{C}}$ can be interpreted as a
well-determined rational function
\[ f(z) = P(z)/Q(z) \in \mathbb{C}(z) \]
the quotient of two well-determined relatively prime polynomials,
with $Q$ monic, satisfying algebraic conditions
that express, in particular, the fact that $f$ is
unramified outside the values $0$, $1$, $\infty$, and imply that the
coefficients of these polynomials are \uline{algebraic numbers}; thus
their zeros are algebraic numbers, representing respectively
the vertices and the centers of the faces of the map under
consideration.

Returning to the general case, since finite maps are interpreted
as coverings over $\bar{\mathbb{Q}}$ of an algebraic curve defined
over the prime field $\mathbb{Q}$ itself, it follows that the Galois
group $\mathrm{I}\!\Gamma$ of $\bar{\mathbb{Q}}$ over $\mathbb{Q}$ acts naturally on the category of these
maps. For example, the action of an automorphism
$\gamma \in \mathrm{I}\!\Gamma$ on a spherical map given by the rational function
above is obtained by applying $\gamma$ to the coefficients of the
polynomials $P$, $Q$. Here, then, is this mysterious group $\mathrm{I}\!\Gamma$ intervening as
a transforming agent on topological-combinatorial shapes of the
most elementary kind imaginable, prompting questions
such as: are these given oriented maps ``conjugate'',
or: what exactly are the conjugates\frapparatus{V-FR-015-001} of this given oriented
map ? (There are, visibly, only finitely many.)

I have dealt with a few concrete cases (for coverings
of low degrees) by various expedients; J. Malgoire has dealt with
a few others - I doubt that there is a uniform method of
answering them by computer calculations. My reflection very quickly

\margin{15}
\copytwo{A2-119-091-001}{\scancontext{https://grothendieck.umontpellier.fr/web/viewer.html?file=../119.pdf\#page=92}{\href{https://grothendieck.umontpellier.fr/web/viewer.html?file=../119.pdf\#page=92}{[copy 2, scan 119/91]} Copy 2 has the handwritten foliation 91 at the lower left.}}

\finishscan
% END PAGE esq_119_015
% BEGIN PAGE esq_119_016
\scanpage{119/16}
\mspage{14}

set off in a more conceptual direction\frapparatus{V-FR-016-001}, to come to
grasp the nature of this action of $\mathrm{I\!\Gamma}$. One immediately sees
that, roughly speaking, this action is expressed by a certain
``outer'' action of $\mathrm{I\!\Gamma}$ on the profinite compactification\tnote{TN-016-01}{``Profinite compactification'' renders ``compactifié profini'', the expression used here. The source used ``complétions profinies'' on p.~6, translated there as ``profinite completions''.}
of the oriented cartographic group $\underline{C}_2^{+}$ , and this action in turn
is deduced, by passage to the quotient, from the canonical outer
action of $\mathrm{I\!\Gamma}$ on the profinite fundamental group $\hat{\pi}_{0,3}$ of $(U_{0,3})_{\bar{\mathbb{Q}}}$ ,
where $U_{0,3}$ denotes the standard curve of genus $0$ over the prime field
$\mathbb{Q}$ with three points removed. This is how my attention
turned toward what I have since called ``\uline{anabelian}
\uline{algebraic geometry}'', whose starting point is precisely a study
(for the moment restricted to characteristic zero) of the action
of ``absolute'' Galois groups (\unsure{in particular}\tnote{TN-016-02}{The French alternatives are ``nctamment|notamment''. Both are represented by ``in particular''; the letterform doubt remains linked to the source.}
the groups $\mathrm{Gal}(\bar{K}/K)$,
where $K$ is an extension of finite type of the prime field) on
geometric (profinite) fundamental groups of algebraic varieties
(defined over $K$), and more particularly (breaking with
a deeply rooted tradition) fundamental groups that are
very far from abelian groups (and that I therefore
call ``\uline{anabelian}''). Among these groups, and very close to the group
$\hat{\pi}_{0,3}$ , is the profinite compactification of the modular group $S\ell(2,\mathbb{Z})$,
whose quotient by the center $\pm 1$ contains the preceding group as
the congruence subgroup mod 2, and can moreover also be interpreted
as an oriented ``cartographic'' group, namely the one
that classifies \uline{triangulated} oriented maps (i.e. those whose
faces are triangles or monogons).

Every finite oriented map gives rise to a smooth projective
algebraic curve defined over $\bar{\mathbb{Q}}$ , and the question then immediately
arises: which algebraic curves over $\bar{\mathbb{Q}}$ are obtained
in this way - might one obtain them all, who knows ?
In more learned terms, might it be true that every smooth projective
algebraic curve defined over a number field would occur
as a possible ``modular \unsure{curve}''\tnote{TN-016-03}{The French reading remains ``ccurbe|courbe''. The English keeps the doubt on ``curve'' and supplies both source candidates here.}
parametrizing elliptic curves
equipped with a suitable rigidification ? Such a
supposition looked so crazy that I was almost

\margin{16}
\copytwo{A2-119-092-001}{\scancontext{https://grothendieck.umontpellier.fr/web/viewer.html?file=../119.pdf\#page=93}{\href{https://grothendieck.umontpellier.fr/web/viewer.html?file=../119.pdf\#page=93}{[copy 2, scan 119/92]} Copy 2 has the handwritten foliation 92 at the lower left.}}

\finishscan
% END PAGE esq_119_016
% BEGIN PAGE esq_119_017
\scanpage{119/17}
\mspage{15}

embarrassed to submit it to the experts in the subject. Deligne, consulted,
did indeed find the supposition crazy\frapparatus{V-FR-017-001}, but had no counterexample
up his sleeve. Less than a year later, at the International
Congress in Helsinki, the Soviet mathematician Bielyi\tnote{TN-017-01}{The source's spelling ``Bielyi'' is retained.}
announced precisely this result, with a proof of disconcerting simplicity
occupying two short pages of a letter from Deligne - probably never
has a deep and bewildering result been proved
in so few lines !

In the form in which Bielyi states it, his result says essentially
that \uline{every algebraic curve defined over a number}
\uline{field can be obtained as a covering of the projective line ramified}
\uline{only at the points} 0, 1, $\infty$. This result seems to have passed
more or less unnoticed. Yet it appears to me to have considerable
scope. For me, its essential message was that \uline{there is a deep}
\uline{identity between the combinatorics of finite maps on the one hand}, \uline{and the}
\uline{geometry of algebraic curves defined over number fields}
\uline{on the other}. This deep result, together with the algebraic-geometric
interpretation of finite maps, opens the door onto a new,
unexplored world - and within reach of everyone, who passes by without seeing it.

It was only almost three years later, seeing that the vast
horizons opening up here decidedly stirred
nothing in any of my students, or even in any of the three or
four colleagues of the highest caliber to whom I had the opportunity
to speak about it in detail, that I made a first exploratory
journey into this ``new world'', from January to June 1981. This first venture
took material form in a packet of some 1 300 manuscript pages,
christened ``The Long March through Galois Theory''.\tnote{TN-017-02}{The French title is ``La Longue Marche à travers la théorie de Galois''.}
It is above all an effort to understand the relations between
``arithmetic'' Galois groups and ``geometric'' profinite fundamental
groups. Fairly quickly it turned toward formulating
the action of $\mathrm{Gal}(\bar{\mathbb{Q}}/\mathbb{Q})$ on $\hat{\pi}_{0,3}$ in calculational terms, and at a
later stage, on the slightly larger group $S\ell(2,\mathbb{Z})^{\wedge}$, which
gives rise to a more elegant and more effective formalism. It was also
in the course of this work (but developed in separate notes)

\margin{17}
\copytwo{A2-119-093-001}{\scancontext{https://grothendieck.umontpellier.fr/web/viewer.html?file=../119.pdf\#page=94}{\href{https://grothendieck.umontpellier.fr/web/viewer.html?file=../119.pdf\#page=94}{[copy 2, scan 119/93]} Copy 2 has the handwritten foliation 93 at the lower left.}}

\finishscan
% END PAGE esq_119_017
% BEGIN PAGE esq_119_018
\scanpage{119/18}
\mspage{16}

that the central theme of anabelian algebraic geometry appeared:
reconstructing certain varieties $X$ called ``anabelian''
over an absolute field $K$ from their mixed fundamental group,
an extension of $\mathrm{Gal}(\bar{K}/K)$ by $\pi_1(X_{\bar{K}})$; it was then that the
``fundamental conjecture of anabelian algebraic geometry'' emerged,
close to the conjectures of Mordell and Tate that Faltings
has just proved\notecall{copy1-3}{call-esq_119_018-3-1}{($^3$)}. It was also there that a first reflection on
Teichmüller groups and the first intuitions about the multiple
structure of the ``Teichmüller tower'' began - the open moduli
multiplicities $M_{g,\nu}$ appearing, moreover, as the first important
examples, in dimension $> 1$, of varieties (or rather,
multiplicities) that seem well deserving of the name ``anabelian''.
Toward the end of this period of reflection, it appeared to me
as a foundational reflection on a theory still in
limbo, for which the name ``\uline{Galois-Teichmüller theory}''
seemed more appropriate than ``Galois Theory'', the name I had initially
given to my notes.

This is not the place to give a more detailed account
of this whole collection of questions, intuitions, ideas - including
tangible results, certainly. The most important seems to me to be the one
mentioned in passing in par. 2, namely the faithfulness of the outer
action of $\mathrm{I}\!\Gamma = \mathrm{Gal}(\bar{\mathbb{Q}}/\mathbb{Q})$ (and of its open subgroups) on $\hat{\pi}_{0,3}$,
and more generally (if I remember correctly) on the fundamental
group of any ``anabelian'' algebraic curve (i.e. whose genus
$g$ and ``number of holes'' $\nu$ satisfy the inequality $2g + \nu \geq 3$,
i.e. such that $\chi(X) < 0$) defined over a finite extension of $\mathbb{Q}$. This
result can be regarded as essentially equivalent to
Bielyi's theorem - it is the first concrete manifestation, through
a precise mathematical statement, of the ``message'' discussed
above.

I would like to finish this rapid overview with a few words
of commentary on the truly unimaginable richness of a typical
anabelian group such as the group $S\ell(2,\mathbb{Z})$ - probably the most
remarkable infinite discrete group one has encountered, which appears

\margin{18}
\copytwo{A2-119-094-001}{\scancontext{https://grothendieck.umontpellier.fr/web/viewer.html?file=../119.pdf\#page=95}{\href{https://grothendieck.umontpellier.fr/web/viewer.html?file=../119.pdf\#page=95}{[copy 2, scan 119/94]} Copy 2 has the handwritten foliation 94 at the lower left.}}

\finishscan
% END PAGE esq_119_018
% BEGIN PAGE esq_119_019
\scanpage{119/19}
\mspage{17}

in a multitude of guises (some of which have been touched on
in the present report), and which, from the viewpoint of
Galois-Teichmüller theory, can be regarded as the fundamental
``building block'' of the ``Teichmüller tower''. The structural
element of $S\ell(2,\mathbb{Z})$ that fascinates me above all is of course
the outer action of the Galois group $\mathrm{I}\!\Gamma$ on its profinite
compactification. By Bielyi's theorem, taking the profinite
compactifications of subgroups of finite index of $S\ell(2,\mathbb{Z})$, and the induced
outer action (passing also to an open subgroup of $\mathrm{I}\!\Gamma$
as necessary), \uline{one essentially obtains the fundamental groups}
\uline{of all algebraic curves} (not necessarily compact)
\uline{defined over number fields} $K$, \uline{and the outer action of}
$\mathrm{Gal}(\bar{K}/K)$ \uline{on them} - at least it is true that every such fundamental
group appears as a quotient of one of the former groups (*).
Taking into account the ``anabelian yoga'' (which remains conjectural),
according to which an anabelian algebraic curve over a number field $K$
(a finite extension of $\mathbb{Q}$) is known up to isomorphism\frapparatus{V-FR-019-001} when
its mixed fundamental group is known (or, what amounts to the same thing,
the outer action of $\mathrm{Gal}(\bar{K}/K)$ on its geometric profinite
fundamental group), one can thus say that \uline{all algebraic curves}
\uline{defined over number fields are} ``\uline{contained}'' \uline{in the profinite}
\uline{compactification} $S\ell(2,\mathbb{Z})\hat{}$, \uline{and in the knowledge of a certain}
\uline{subgroup} $\mathrm{I}\!\Gamma$ \uline{of its group of outer automorphisms} !
Passing to the abelianizations of the preceding fundamental groups, one sees,
in particular, that all the abelian $\ell$-adic representations dear
to Tate and his associates, defined by Jacobians and generalized
Jacobians of algebraic curves defined over number fields,
are contained in this single action of $\mathrm{I}\!\Gamma$ on the anabelian profinite
group $S\ell(2,\mathbb{Z})\hat{}$!\notecall{copy1-4}{call-esq_119_019-4-1}{($^{4}$)}

\rule{5cm}{0.4pt}

(*) In fact, these are quotients of a particularly trivial
kind, by abelian subgroups that are products of ``Tate modules''
$\hat{\mathbb{Z}}(1)$, corresponding to ``groups of loops'' about points at infinity.\tnote{TN-019-01}{The expression is ``groupes-lacets'' in the French; the wording ``groups of loops'' preserves the author's phrase at this place.}

\margin{19}
\copytwo{A2-119-095-001}{\scancontext{https://grothendieck.umontpellier.fr/web/viewer.html?file=../119.pdf\#page=96}{\href{https://grothendieck.umontpellier.fr/web/viewer.html?file=../119.pdf\#page=96}{[copy 2, scan 119/95]} Copy 2 has a handwritten question mark in the left margin beside “damental apparaît comme quotient”.}}
\copytwo{A2-119-095-002}{\scancontext{https://grothendieck.umontpellier.fr/web/viewer.html?file=../119.pdf\#page=96}{\href{https://grothendieck.umontpellier.fr/web/viewer.html?file=../119.pdf\#page=96}{[copy 2, scan 119/95]} Copy 2 has a handwritten oval around the typed footnote call “(*).”, including its following full stop.}}
\copytwo{A2-119-095-003}{\scancontext{https://grothendieck.umontpellier.fr/web/viewer.html?file=../119.pdf\#page=96}{\href{https://grothendieck.umontpellier.fr/web/viewer.html?file=../119.pdf\#page=96}{[copy 2, scan 119/95]} Copy 2 has another handwritten question mark in the left margin beside “viale, par des sous-groupes” in the footnote.}}
\copytwo{A2-119-095-004}{\scancontext{https://grothendieck.umontpellier.fr/web/viewer.html?file=../119.pdf\#page=96}{\href{https://grothendieck.umontpellier.fr/web/viewer.html?file=../119.pdf\#page=96}{[copy 2, scan 119/95]} Copy 2 has the handwritten foliation 95 at the lower left.}}

\finishscan
% END PAGE esq_119_019
% BEGIN PAGE esq_119_020
\scanpage{119/20}
\mspage{18}

Faced with this, some are content to shrug their shoulders
with a disillusioned air and bet that nothing can be drawn
from it all except dreams. They forget, or are unaware,
that our science, and every science, would amount to very little
if, since its beginnings, it had not been nourished by the dreams and visions
of those who devote themselves to it with passion.

4.        From the beginning of my reflection on two-dimensional maps,
I was especially interested in the so-called
``\uline{regular}'' maps, that is, those whose automorphism group
acts transitively (and consequently, simply transitively)
on the set of frames. In the oriented case and in terms of
the algebraic-geometric interpretation in the preceding paragraph,
these are the maps that correspond to a \uline{Galois} covering of
the projective line. Very quickly, too, and even before the connection
with algebraic geometry appeared, it also became
necessary not to exclude infinite maps, which in particular enter
naturally as universal coverings of
finite maps. It appears (as an immediate consequence of the ``dictionary''
of maps, extended to maps that need not be finite)
that for every pair of natural numbers $p, q$ $\geq$ 1, there exists up to
(nonunique) isomorphism one and only one 1-\uline{connected} map
of type $(p,q)$, i.e. all of whose vertices have order $p$
and all of whose faces have order $q$, and this map is regular.
It is pinned down by the choice of a frame, and its
automorphism group is then canonically isomorphic to the quotient
of the cartographic group (resp. of the oriented cartographic
group in the oriented case) by the additional relations
\[ \rho_s^{\,p} = \rho_f^{\,q} = 1. \]
The case in which this group is finite is the ``Pythagorean'' case of
regular spherical maps; the case in which it is infinite gives the regular tessellations

\margin{20}
\copytwo{A2-119-096-001}{\scancontext{https://grothendieck.umontpellier.fr/web/viewer.html?file=../119.pdf\#page=97}{\href{https://grothendieck.umontpellier.fr/web/viewer.html?file=../119.pdf\#page=97}{[copy 2, scan 119/96]} Copy 2 has the handwritten foliation 96 at the lower left.}}

\finishscan
% END PAGE esq_119_020
% BEGIN PAGE esq_119_021
\scanpage{119/21}
\mspage{19}

of the euclidean or hyperbolic plane (*). The connection
of the combinatorial theory with the ``conformal'' theory of regular
tessellations of the hyperbolic plane was sensed before
the connection of finite maps with finite coverings of the
projective line appeared. Once this connection is understood, it becomes obvious
that it must also extend to infinite maps (regular
or otherwise): \uline{every map, finite or otherwise, is canonically realized on
a conformal surface} (compact if and only if the map is
finite), \uline{as a ramified covering of the complex projective line,
ramified only at the points} 0, 1, $\infty$. The only difficulty
here was to work out the dictionary between topological
maps and sets with operators, which posed some conceptual
problems in the infinite case, beginning with the very notion of
a ``topological map''. In particular, both for the
internal coherence of the dictionary and so as not to miss
certain interesting cases of infinite maps, it proves necessary
not to exclude vertices and faces of infinite order. This foundational
work was also carried out by J. Malgoire and C. Voisin, following
on from their first work on finite maps, and their theory
indeed supplies everything one was entitled to expect from it
(and even more...).

It was in 1977 and 1978, alongside two C 4 courses\tnote{TN-021-01}{``C 4'' is the course designation used in the source; it is retained here.}
on the geometry of the cube and that of the icosahedron, that I
began to take an interest in regular polyhedra, which then appeared
to me as particularly concrete ``geometric realizations''
of combinatorial maps, the vertices, edges and faces being
realized respectively as points, lines and planes
in a suitable three-dimensional affine space, with the incidence
relations respected. This notion of geometric realization

\rule{4cm}{0.4pt}

(*) In these statements, one should not exclude the case in which $p$, $q$
can take the value $+\infty$, which one encounters in particular very
naturally in tessellations associated with certain infinite regular
polyhedra, cf. below.

\margin{21}
\copytwo{A2-119-097-001}{\scancontext{https://grothendieck.umontpellier.fr/web/viewer.html?file=../119.pdf\#page=98}{\href{https://grothendieck.umontpellier.fr/web/viewer.html?file=../119.pdf\#page=98}{[copy 2, scan 119/97]} Copy 2 has the handwritten foliation 97 at the lower left.}}

\finishscan
% END PAGE esq_119_021
% BEGIN PAGE esq_119_022
\scanpage{119/22}
\mspage{20}

of a combinatorial map retains a meaning over a base field, and
even over an arbitrary base ring. It also retains a
meaning for regular polyhedra of any dimension, by
replacing the cartographic group $\underline{C}_2$ with a suitable n-dimensional
variant $\underline{C}_n$. The case n = 1, i.e. the theory of
regular polygons in any characteristic, was the subject
of a DEA course in 1977/78, and already brought out some
new phenomena, as well as the usefulness of working in a \uline{projective}
space rather than an affine ambient space (here the affine plane).
This is due in particular to the fact that in certain
characteristics (and especially in characteristic 2) the
center of a regular polyhedron is thrust out to infinity. On the other
hand, the projective context, unlike the affine context,
allows one to develop with ease a duality formalism for
regular polyhedra, corresponding to the duality formalism
of combinatorial or topological maps (where the roles of vertices
and faces, in the case n = 2, say, are interchanged).
It turns out that for every projective regular polyhedron, one can
define an associated canonical hyperplane, which plays the role of
a canonical hyperplane at infinity and allows one to regard the given
polyhedron as an affine regular polyhedron.

The extension of the theory of regular polyhedra (and,
more generally, of all sorts of geometric-combinatorial
configurations, including root systems...) from the base field
$\mathbb{R}$ or $\mathbb{C}$ to a general base ring seems to me to have,
in this part of geometry, a scope comparable to the analogous
extension that has taken place since the beginning of the century in algebraic
geometry, or over the last twenty years or so in topology ( *),
with the introduction of the language of schemes and that of topoi.

\rule{4cm}{0.4pt}

(*) In writing this, I am aware that few topologists,
even today, recognize this conceptual and technical
broadening of topology and the resources
it offers.

\margin{22}
\copytwo{A2-119-098-001}{\scancontext{https://grothendieck.umontpellier.fr/web/viewer.html?file=../119.pdf\#page=99}{\href{https://grothendieck.umontpellier.fr/web/viewer.html?file=../119.pdf\#page=99}{[copy 2, scan 119/98]} Copy 2 has the handwritten foliation 98 at the lower left.}}

\finishscan
% END PAGE esq_119_022
% BEGIN PAGE esq_119_023
\scanpage{119/23}
\mspage{21}

My sporadic reflection on this question over several years
was confined to bringing out a few simple basic principles, at
first directing my attention to \uline{pinned} regular
polyhedra,\tnote{TN-023-01}{``Pinned'' renders ``épinglé'', continuing the source's language of fixing a polyhedron by choosing a frame.}
which reduces the necessary conceptual baggage to a minimum
and practically eliminates rationality questions that are at all
delicate. For such a polyhedron one finds a canonical basis (or
frame) of the ambient affine or projective space, in
such a way that the operations of the cartographic group $\underline{C}_n$, generated
by the fundamental reflections $\sigma_i$ $(0 \leq i \leq n)$, are written in it
by universal formulas in terms of $n$ parameters
$\alpha_1,\ldots,\alpha_n$ , which geometrically are interpreted as twice
the cosines of the ``fundamental angles'' of the polyhedron. The polyhedron
is reconstructed from this action and the affine or
projective flag associated with the chosen basis, by transforming that flag
by all the elements of the group generated by the fundamental
reflections. Thus the ``universal'' pinned n-polyhedron is
canonically defined over the polynomial ring in $n$ indeterminates
\[ \mathbb{Z}[\underline{\alpha}_1,\ldots,\underline{\alpha}_n] , \]
its specializations over arbitrary base fields $k$ (through
values $\alpha_i \in k$ assigned to the indeterminates $\underline{\alpha}_i$) giving
regular polyhedra corresponding to various combinatorial
types. In this game there is no question of restricting oneself to
finite regular polyhedra, or even to regular polyhedra whose
facets have finite order, i.e. for which the parameters
$\alpha_i$ are roots of suitable ``semicyclotomic'' equations
expressing that the ``fundamental angles'' (in the
case where the base field is $\mathbb{R}$) are commensurable with $2\pi$. Already
when $n = 1$, perhaps the most interesting regular polygon
of all (morally, the regular polygon with just one side !)
is the one corresponding to $\alpha = 2$, giving rise to a circumscribed
parabolic conic, i.e. tangent to the line at infinity.

The finite case is the one in which the group generated by the fundamental
reflections, which is also the automorphism group of the

\margin{23}
\copytwo{A2-119-099-001}{\scancontext{https://grothendieck.umontpellier.fr/web/viewer.html?file=../119.pdf\#page=100}{\href{https://grothendieck.umontpellier.fr/web/viewer.html?file=../119.pdf\#page=100}{[copy 2, scan 119/99]} Copy 2 has the handwritten foliation 99 at the lower left.}}

\finishscan
% END PAGE esq_119_023
% BEGIN PAGE esq_119_024
\scanpage{119/24}
\mspage{22}

regular polyhedron under consideration is finite. In the case of the
base field $\mathbb{R}$ (or $\mathbb{C}$, which amounts to the same thing), and for $n = 2$, the finite
cases have been well known since antiquity - which does not exclude
the scheme-theoretic viewpoint's revealing new charms\frapparatus{V-FR-024-001}
in them; one can say, however, that by specializing the icosahedron\frapparatus{V-FR-024-002}
(for example) over finite base fields of arbitrary
characteristic, one always obtains an icosahedron, with its own
combinatorics and the same simple automorphism group of order 60.
The same remark applies to finite regular polyhedra
of higher dimension, studied systematically in two
fine books by Coxeter. The situation is quite different on starting
with an \uline{infinite} regular polyhedron, over a field such as $\mathbb{Q}$, say,
and ``specializing'' it over the prime fields $\mathbb{F}_p$ (an operation
well defined for every $p$ except finitely many primes).
It is clear that every regular polyhedron over a finite field is finite -
\uline{one thus finds infinitely many finite regular polyhedra for} $p$
\uline{varying}, \uline{whose combinatorial type}, or, what amounts to the same thing,
whose automorphism group, \uline{varies} ``\uline{arithmetically}'' with p.
This situation is particularly intriguing in the case
$n = 2$, where one has the relation made explicit in the preceding
paragraph between combinatorial 2-maps and algebraic curves
defined over number fields. In this case an infinite regular
polyhedron defined over any infinite field (and hence
over a sub-$\mathbb{Z}$-algebra with two generators \unsure{.}of it)\tnote{TN-024-01}{The small dot before the French ``de'' is unresolved; it remains at the corresponding ``of'' in the translation.}
thus gives rise to infinitely many algebraic curves defined over
number fields, which are Galois coverings of the standard projective
line ramified only at 0, 1, $\infty$. The optimal case is of course
that in which one starts with the universal regular 2-polyhedron,
or rather with the one deduced from it by passing to the fraction
field $\mathbb{Q}(\alpha_1,\alpha_2)$ of its base ring. This raises a host of
new questions, vague as well as precise, none of which
I have yet had the leisure to examine more closely - I shall
mention only this one: what exactly are the finite regular
2-maps, or, what amounts to the same thing, \unsure{.}the finite quotient groups\tnote{TN-024-02}{The unresolved dot before the French ``les'' is retained before the corresponding English ``the''.}

\margin{24}
\copytwo{A2-119-100-001}{\scancontext{https://grothendieck.umontpellier.fr/web/viewer.html?file=../119.pdf\#page=101}{\href{https://grothendieck.umontpellier.fr/web/viewer.html?file=../119.pdf\#page=101}{[copy 2, scan 119/100]} Copy 2 has the handwritten foliation 100 at the lower left.}}

\finishscan
% END PAGE esq_119_024
% BEGIN PAGE esq_119_025
\scanpage{119/25}
\mspage{23}

of the 2-cartographic group that arise from regular 2-polyhedra
over finite fields (*) ? Would one obtain them all,
and, if so, how ?

These reflections bring into full light the fact,
which for me was entirely unexpected, that the theory of
finite regular polyhedra, already in the case of dimension
$n = 2$, is infinitely richer, and in particular gives infinitely
more different combinatorial forms, when one admits
base fields of nonzero characteristic,\quad than in the
case considered up to now, where the base fields were restricted
to $\mathbb{R}$ ,or, at a stretch, $\mathbb{C}$ (in the case of what Coxeter
calls ``complex regular polyhedra'', which I prefer
to call ``regular pseudo-polyhedra defined over $\mathbb{C}$ '') (**). Moreover,
it seems that this broadening of viewpoint should also cast
new light on the cases already known. Thus, examining the
Pythagorean polyhedra one after another, I saw
the same little miracle repeated each time, which I called the \uline{combinatorial
paradigm} of the polyhedron in question. Speaking vaguely, it can be
described by saying that, on looking at the specialization of the polyhedron
in the characteristic, or one of the characteristics, that is
the most singular one (or ones) (these are characteristics 2 and 5\frapparatus{V-FR-025-001}

\rule{5cm}{0.4pt}

(*) These are, moreover, the same as those arising from regular
polyhedra over any fields, or over algebraically closed fields, as
one sees by standard specialization arguments.

(**) Pinned pseudo-polyhedra are described in the same way as
pinned polyhedra, with the sole difference that the fundamental
reflections $\sigma_i$ $(0 \leq i \leq n)$ are replaced here by \uline{pseudo-reflections}
(which Coxeter moreover assumes to have finite order, as he
restricts himself to finite combinatorial structures). This simply leads
to the introduction, for each of the \unsure{.}$\sigma_i$, of an additional numerical
invariant $\beta_i$ , so that the universal n-pseudo-polyhedron can still be
defined over a polynomial ring with integer coefficients, in
the n+(n+1) variables $\underline{\alpha}_i$ $(1 \leq i \leq n)$ and $\underline{\beta}_j$ $(0 \leq j \leq n)$.

\margin{25}
\copytwo{A2-119-101-001}{\scancontext{https://grothendieck.umontpellier.fr/web/viewer.html?file=../119.pdf\#page=102}{\href{https://grothendieck.umontpellier.fr/web/viewer.html?file=../119.pdf\#page=102}{[copy 2, scan 119/101]} Copy 2 has the handwritten foliation 101 at the lower left.}}

\finishscan
% END PAGE esq_119_025
% BEGIN PAGE esq_119_026
\scanpage{119/26}
\mspage{24}

for the icosahedron, characteristic 2 for the octahedron), one reads
off the geometric regular polyhedron over the finite field concerned
($\mathbb{F}_4$ and $\mathbb{F}_5$ for the icosahedron, $\mathbb{F}_2$ for the octahedron) a particularly
elegant (and unexpected) description of the combinatorics
of the polyhedron. I even seemed to glimpse here a principle
of great generality, which I believed I found again, in particular, in
a later reflection on the combinatorics of the system of 27
lines on a cubic surface and its relations with the root
system $E_7$. Whether such a principle really exists and one
even succeeds in bringing it out\frapparatus{V-FR-026-001} of its cloak of mist, or whether it
recedes as one pursues it and ends by
vanishing like a Fata Morgana,\tnote{TN-026-01}{``Fata Morgana'' is retained as the author's image of an elusive mirage.}
for my part I \unsure{find}\tnote{TN-026-02}{The French alternatives remain ``trcuve|trouve''; both are represented by ``find'', with the doubt linked to the scan.}
in it an uncommon motivating force and fascination,
like that of a dream, perhaps. There is no doubt that following such a call from the
unformulated, from the formless seeking form, from an elusive
glimpse that seems to take pleasure both in evading one and in
revealing itself can only lead far, while no one could
predict where ...

Yet, occupied with other interests and tasks, I have
not so far followed this call, nor met anyone else
who wished to hear it, still less to follow it. Apart from
a few digressions toward other kinds of geometric-combinatorial
structures, my work here too has been confined to an initial
roughing-out and housekeeping, on which it is unnecessary
to dwell further here\notecall{copy1-5}{call-esq_119_026-5-1}{\hand{$^{(5)}$}}. The only point that perhaps still deserves
mention is the existence and uniqueness of the hyperquadric circumscribed
about a given regular n-polyhedron, whose equation can be made
explicit by simple formulas in terms of the fundamental parameters
$\alpha_i$ (*). The case that interests me most is the one where $n = 2$, and the

\rule{5cm}{0.4pt}

(*) An analogous result holds for pseudo-polyhedra\unsure{.|:} It would seem
that the ``exceptional characteristics'' discussed above,
for the specializations of a given polyhedron, are those for
which the circumscribed \var{hyperquadratic}{hyperquadric}{slip corrected}{V-FR-026-002} is either degenerate or
tangent to the hyperplane at infinity.

\margin{26}
\copytwo{A2-119-102-002}{\scancontext{https://grothendieck.umontpellier.fr/web/viewer.html?file=../119.pdf\#page=103}{\href{https://grothendieck.umontpellier.fr/web/viewer.html?file=../119.pdf\#page=103}{[copy 2, scan 119/102]} Copy 2 has a small handwritten cross in the left margin of the footnote, beside the line beginning “lesquelles”.}}
\copytwo{A2-119-102-003}{\scancontext{https://grothendieck.umontpellier.fr/web/viewer.html?file=../119.pdf\#page=103}{\href{https://grothendieck.umontpellier.fr/web/viewer.html?file=../119.pdf\#page=103}{[copy 2, scan 119/102]} Copy 2 has a period after “pseudo-polyèdres”; copy 1 retains the doubtful reading \unsure{.|:} at this position.}}
\copytwo{A2-119-102-004}{\scancontext{https://grothendieck.umontpellier.fr/web/viewer.html?file=../119.pdf\#page=103}{\href{https://grothendieck.umontpellier.fr/web/viewer.html?file=../119.pdf\#page=103}{[copy 2, scan 119/102]} Copy 2 has handwritten “102” at bottom left.}}

\finishscan
% END PAGE esq_119_026
% BEGIN PAGE esq_119_027
\scanpage{119/27}
\mspage{25}

time seems ripe to rewrite a new version, in a modern
style, of Klein's classical book on the icosahedron and the other
Pythagorean polyhedra. Writing such an account of regular
2-polyhedra would be a splendid opportunity for a young researcher
to become familiar both with the geometry of polyhedra and
their connections with spherical\frapparatus{V-FR-027-002}, euclidean and hyperbolic geometries,
and with algebraic curves, and with the language and
basic techniques of modern algebraic geometry. Will one
ever be found to seize this opportunity ?

5. \qquad I would now like to say a few words about certain
reflections that led me to understand the need for new foundations
for ``geometric'' topology, in a direction quite
different from the notion of a topos, and independent even of the needs
of so-called ``abstract'' algebraic geometry (over general base fields
and rings). The initial problem, which began to
intrigue me what must already be some fifteen years ago, was
that of defining a theory of ``dévissage''\tnote{TN-027-01}{``Dévissage'' is retained, following the terminology standard. The source explains it here through reconstruction from canonical building pieces.}
of stratified structures, in order to reconstruct them by a canonical process
from ``building blocks'' canonically deduced from the given
structure. Probably the principal example that had then led me
to this question was the canonical stratification of a
singular algebraic variety\frapparatus{V-FR-027-001} (or a singular complex or real analytic
space) by the decreasing sequence of its successive
``singular loci''. But I probably already sensed the ubiquity
of stratified structures in practically every domain of
geometry (which others surely saw clearly long before
I did). Since then, I have seen such structures appear, in particular,
in every ``moduli'' situation for geometric objects susceptible
not only to continuous variation but, at the same time,
to phenomena of ``degeneracy'' (or ``specialization'') - the
strata then corresponding to the various ``levels of singularity''
(or to the associated combinatorial types) of the objects considered.

\margin{27}
\copytwo{A2-119-103-001}{\scancontext{https://grothendieck.umontpellier.fr/web/viewer.html?file=../119.pdf\#page=104}{\href{https://grothendieck.umontpellier.fr/web/viewer.html?file=../119.pdf\#page=104}{[copy 2, scan 119/103]} Copy 2 has handwritten “103” at bottom left.}}

\finishscan
% END PAGE esq_119_027
% BEGIN PAGE esq_119_028
\scanpage{119/28}
\mspage{26}

The compactified Mumford-Deligne moduli multiplicities $M^{\wedge}_{g,\nu}$
for stable algebraic curves of type $(g,\nu)$
provide a typical and particularly inspiring example, which
played an important motivating role in the resumption of my reflection
on stratified structures, from December 1981 to January
1982. Two-dimensional geometry provides many other
examples of such stratified moduli structures, all of which,
moreover (apart from expedients of rigidification\frapparatus{V-FR-028-001}), appear
as ``multiplicities'' rather than as spaces or varieties
in the ordinary sense (the points of these multiplicities\frapparatus{V-FR-028-002} can have
nontrivial automorphism groups). Among the objects of
two-dimensional geometry giving rise to such stratified moduli
structures of arbitrary dimension, or even infinite
dimension, I shall mention polygons (euclidean, spherical or
hyperbolic), systems of lines in a plane (projective,
say), systems of ``pseudolines'' in a topological projective
plane, or more general immersed curves with normal \unsure{crossings}\tnote{TN-028-01}{The French alternatives are ``crc isements|croisements''. Both are represented by ``crossings'', with the unresolved letterform still marked and linked.}
in a given surface (compact, say).

The simplest nontrivial example of a stratified
structure is obtained by considering a pair $(X,Y)$ consisting of a space
$X$ and a closed subspace $Y$, making a suitable equisingularity
assumption on $X$ along $Y$, and assuming
further (to fix ideas) that the two strata $Y$ and $X\setminus Y$ are
topological \uline{manifolds}. The naive idea, in such a situation,
is to take ``the'' tubular neighborhood $T$ of $Y$ in $X$, whose
boundary $\partial T$ should also be a smooth manifold, fibered over
$Y$ with smooth compact fibers, $T$ itself being identified\frapparatus{V-FR-028-003} with
the bundle of cones over $T$ associated with the preceding bundle. Setting
\[ U = X \setminus \mathrm{Int}(T), \]
one finds a manifold with boundary whose boundary is canonically
isomorphic to that of $T$. That said, the envisaged ``building blocks''
are the manifold with boundary $U$ (compact when $X$ was compact, and
replacing the ``open'' stratum $X\setminus Y$ while making it more precise) and the manifold

\margin{28}
\copytwo{A2-119-104-002}{\scancontext{https://grothendieck.umontpellier.fr/web/viewer.html?file=../119.pdf\#page=105}{\href{https://grothendieck.umontpellier.fr/web/viewer.html?file=../119.pdf\#page=105}{[copy 2, scan 119/104]} Copy 2 has handwritten “104” at bottom left.}}

\finishscan
% END PAGE esq_119_028
% BEGIN PAGE esq_119_029
\scanpage{119/29}
\mspage{27}

(without boundary) Y, with, as additional structure linking them,
the so-called ``gluing'' map
\[ f : \partial U \longrightarrow Y \]
which is a proper smooth fibration. The initial situation
(X,Y) is reconstructed from $(U,Y, f : \partial U \longrightarrow Y)$ by the
formula
\[ X \simeq U \amalg_{\partial U} Y \]
(amalgamated sum under $\partial U$, mapping to U and Y through the inclusion,
resp. the gluing map).

This naive vision immediately encounters various
difficulties. The first is the somewhat vague nature of the
very notion of a tubular neighborhood, which acquires an at all
precise meaning only in the presence of much more rigid structures
than the topological structure alone, such as the ``piecewise
linear'' or riemannian structure (more generally, that of a space with a
distance function); the trouble here is that in none of the examples
that come spontaneously to mind does one naturally have
a structure of this kind - at best one has an equivalence class
of such structures, allowing one to rigidify the situation
just a little. If one moreover admits that an expedient
has been found for obtaining a tubular neighborhood with the desired
properties, which is also unique modulo an automorphism
(topological, say) of the situation, an automorphism that also
respects the fibered structure supplied by the gluing function,
there remains the difficulty that the choices made are not canonical,
the automorphism in question visibly failing to be unique, whatever
one does to ``normalize'' it. The idea here, to make canonical
what is not so, is to work systematically in
``\uline{isotopy categories}'' associated with the categories of a topological
kind that arise in these questions (such as the category of
admissible pairs (X,Y) and homeomorphisms of such pairs\frapparatus{V-FR-029-001},
etc.), keeping the same objects but taking as ``morphisms''

\margin{29}
\copytwo{A2-119-105-001}{\scancontext{https://grothendieck.umontpellier.fr/web/viewer.html?file=../119.pdf\#page=106}{\href{https://grothendieck.umontpellier.fr/web/viewer.html?file=../119.pdf\#page=106}{[copy 2, scan 119/105]} Copy 2 has handwritten “105” at bottom left.}}

\finishscan
% END PAGE esq_119_029
% BEGIN PAGE esq_119_030
\scanpage{119/30}
\mspage{28}

the isotopy classes (in a sense dictated unambiguously by the
context) of isomorphisms (or even of morphisms more general
than isomorphisms). This idea, taken up successfully
in Yves Ladegaillerie's thesis in particular (cf. the beginning of par. 3),
has served me systematically in all my subsequent reflections
on combinatorial topology, whenever it was a matter of formulating
precise theorems translating topological situations
into combinatorial ones. In the present situation,
my hope was to reach the formulation (and proof !)
of an equivalence theorem between two suitable isotopy categories,
one being the category of ``admissible pairs'' $(X,Y)$,
the other that of ``admissible triples'' $(U,Y,f)$, where $Y$ is a
manifold, $U$ a manifold with boundary, and $f : \partial U \longrightarrow Y$ a proper
smooth fibration. Moreover, of course, I hoped that such a statement,
modulo work of an essentially algebraic nature, would
extend of itself to a more sophisticated statement applying to
general stratified structures.

Very quickly it became apparent that there could be no question
of obtaining so ambitious a statement in the context of
topological spaces, because of the perpetual phenomena of
``wildness''. Already when $X$ itself is a manifold and $Y$ is
reduced to a point, one runs up against the difficulty that the cone over
a compact space $Z$ can be a manifold at its vertex without
$Z$ being homeomorphic to a sphere, or even being a manifold. It
was also clear that the contexts of more rigid structures
existing at the time, such as the ``piecewise linear'' context,
were likewise inadequate - one of the prohibitive reasons they
shared being that, for a pair $(U,S)$ consisting of a
``space'' $U$ and a closed subspace $S$, and a gluing
map $f : S \longrightarrow T$, they did not allow one to construct the corresponding
amalgamated sum. It was a few years later that I was informed
of Hironaka's theory of the sets he calls, I believe,
``semianalytic'' (real), which satisfy some of the essential
stability conditions (probably even all of them)

\margin{30}
\copytwo{A2-119-106-001}{\scancontext{https://grothendieck.umontpellier.fr/web/viewer.html?file=../119.pdf\#page=107}{\href{https://grothendieck.umontpellier.fr/web/viewer.html?file=../119.pdf\#page=107}{[copy 2, scan 119/106]} Copy 2 has a small baseline point before “essentielles”: \unsure{.}. Its origin and punctuation status remain unresolved.}}
\copytwo{A2-119-106-002}{\scancontext{https://grothendieck.umontpellier.fr/web/viewer.html?file=../119.pdf\#page=107}{\href{https://grothendieck.umontpellier.fr/web/viewer.html?file=../119.pdf\#page=107}{[copy 2, scan 119/106]} Copy 2 has handwritten “106” at bottom left.}}

\finishscan
% END PAGE esq_119_030
% BEGIN PAGE esq_119_031
\scanpage{119/31}
\mspage{29}

needed to develop a usable context of ``tame
topology''. This immediately restarts a reflection on the foundations
of such a topology, whose need becomes increasingly clear
to me.

With some ten years' hindsight, I would now say,
on this subject, that ``\uline{general topology}'' \uline{was developed}
(in the thirties and forties) \uline{by analysts and for}
\uline{the needs of analysis}, not for the needs of topology
properly speaking, that is, the study of the \uline{topological properties}
\uline{of various geometric shapes}. This inadequacy of the foundations
of topology manifests itself from the outset through
``false problems'' (at least from the viewpoint of the topological
intuition of shapes), such as ``invariance of domain'',
even though Brouwer's solution of the latter leads him to
introduce important new geometric ideas. Even today,
as in the heroic times when one saw for the first time, with anxiety,
curves cheerfully filling squares and cubes,
when one proposes to do topological geometry in the technical
context of topological spaces, one meets parasitic difficulties at
every step arising from wild phenomena.
Thus, outside cases of (very) low dimension, it can
hardly be possible, for a given space $X$ (a compact manifold,
say), to study the homotopy type (say) of the automorphism
group of $X$, or of the space of embeddings, or immersions, etc.,
of $X$ into some other space $Y$ - although one feels that these
invariants should belong to the arsenal of essential invariants associated
with $X$, or with the pair $(X,Y)$, etc., on the same footing
as the function space $\underline{\mathrm{Hom}}\,(X,Y)$ familiar in homotopy
topology. Topologists evade the difficulty without confronting it,
falling back on contexts close to the topological context and
less marked by wildness, such as differentiable manifolds,
PL (piecewise linear) spaces, etc., visibly none of which is ``good'',
i.e. stable under the most obvious topological operations,

\margin{31}
\copytwo{A2-119-107-001}{\scancontext{https://grothendieck.umontpellier.fr/web/viewer.html?file=../119.pdf\#page=108}{\href{https://grothendieck.umontpellier.fr/web/viewer.html?file=../119.pdf\#page=108}{[copy 2, scan 119/107]} Copy 2 has handwritten “107” at bottom left.}}

\finishscan
% END PAGE esq_119_031
% BEGIN PAGE esq_119_032
\scanpage{119/32}
\mspage{30}

such as the operations of contraction and gluing (without even proceeding
to operations of the type $X \longrightarrow \mathrm{Aut}(X)$, which take one out of the
paradise of finite-dimensional ``spaces''). That is a way of
beating around the bush ! This situation, as so many times already in
the history of our science, simply brings to light the almost
insurmountable inertia of the mind, burdened by conditioning
of considerable weight, when it comes to looking at a
foundational question, and therefore at the very context in which one lives,
breathes and works - instead of accepting it as an immutable
given. It is surely because of this inertia that it took
millennia for an idea or a reality as childlike as
zero, a group or a topological shape to gain citizenship
in mathematics. It is also surely because of it that the straitjacket
of general topology continues to be patiently dragged along by
generations of topologists, ``wildness'' being borne as
an inescapable fate supposedly rooted in the very nature
of things.

My approach toward possible foundations for a tame
topology was an axiomatic one. Rather than declaring
(which would certainly be perfectly reasonable) that
the ``tame spaces'' sought are precisely, say,
Hironaka's semianalytic spaces\frapparatus{V-FR-032-001}, and then developing in
this context the arsenal of constructions and notions familiar in
topology, together, certainly, with those that until now had been
unable to be developed, and for good reason, I preferred to try to bring out
which of the geometric properties of the notion of a semianalytic
set in a space $\mathbb{R}^{n}$ allow such sets to be used as
local ``models'' for a notion of ``\uline{tame space}'' (in this case,
semianalytic), and what (one hopes !) makes this notion of a
tame space sufficiently flexible actually to serve
as a basic notion for a ``tame topology'' capable of expressing
the topological intuition of shapes with ease. Thus, once
the necessary foundational work has been accomplished, there will appear not
\uline{one} ``tame theory'', but a vast infinitude, ranging from the most

\margin{32}
\copytwo{A2-119-108-001}{\scancontext{https://grothendieck.umontpellier.fr/web/viewer.html?file=../119.pdf\#page=109}{\href{https://grothendieck.umontpellier.fr/web/viewer.html?file=../119.pdf\#page=109}{[copy 2, scan 119/108]} Copy 2 has handwritten “108” at bottom left.}}

\finishscan
% END PAGE esq_119_032
% BEGIN PAGE esq_119_033
\scanpage{119/33}
\mspage{31}

restrictive of all, that of ``piecewise $\bar{\mathbb{Q}}_r$-algebraic spaces''\frapparatus{V-FR-033-001}
(where $\bar{\mathbb{Q}}_r = \bar{\mathbb{Q}} \cap \mathbb{R}$), to the one that (rightly or wrongly) appears to me
probably the broadest, namely that of ``piecewise real analytic
spaces'' (or semianalytic in Hironaka's terminology).
Among the foundational theorems envisaged
in my program is a \uline{comparison theorem} which, speaking
vaguely, will say that \uline{one finds essentially the same isotopy
categories} (or even $\infty$-isotopy categories), whatever
tame theory one works with\notecall{copy1-6}{call-esq_119_033-6-1}{\hand{$^{(6)}$}}. More precisely,
it is a matter of putting one's finger on a system of axioms sufficiently
rich to imply (among many other things !) that, given
two tame theories \hand{$\mathcal{T}$}, \hand{$\mathcal{T}$}' with \hand{$\mathcal{T}$} finer than \hand{$\mathcal{T}$}' (in
an obvious sense), and two T-tame spaces X, Y, which
therefore also define corresponding T'-tame spaces,
the canonical map
\[ \underline{\mathrm{Isom}}_{\hand{\mathcal{T}}}(X,Y) \longrightarrow \underline{\mathrm{Isom}}_{\hand{\mathcal{T}}'}(X,Y) \]
induces a bijection on the set of connected components (which
will imply that the isotopy category of T-spaces is equivalent
to that of T'-spaces), and is even a homotopy
equivalence (which means that one even has an equivalence for
the ``$\infty$-isotopy'' categories\frapparatus{V-FR-033-002}, finer than the isotopy categories
in which only the $\pi_0$ of the isomorphism spaces is retained).
Here the \uline{Isom} can be defined in an obvious way as
semi-simplicial sets, for example, to give a precise meaning
to the preceding statement. Analogous statements should
hold on replacing the \uline{Isom} ``spaces'' by other spaces of
maps subject to \unsure{.}standard geometric conditions, such as
being embeddings, immersions, smooth, étale,
fibrations, etc. One likewise expects analogous
statements in which X, Y are replaced by systems of tame spaces,
such as those entering a theory of dévissage of
stratified structures - so that, in a precise technical
sense, this dévissage theory too will be essentially

\margin{33}
\copytwo{A2-119-109-001}{\scancontext{https://grothendieck.umontpellier.fr/web/viewer.html?file=../119.pdf\#page=110}{\href{https://grothendieck.umontpellier.fr/web/viewer.html?file=../119.pdf\#page=110}{[copy 2, scan 119/109]} Copy 2 has a small baseline point after “elle aussi,”: \unsure{.}. Its origin and punctuation status remain unresolved.}}
\copytwo{A2-119-109-002}{\scancontext{https://grothendieck.umontpellier.fr/web/viewer.html?file=../119.pdf\#page=110}{\href{https://grothendieck.umontpellier.fr/web/viewer.html?file=../119.pdf\#page=110}{[copy 2, scan 119/109]} Copy 2 has handwritten “109” at bottom left.}}

\finishscan
% END PAGE esq_119_033
% BEGIN PAGE esq_119_034
\scanpage{119/34}
\mspage{32}

independent of the tame theory chosen to express
it.

The first decisive test for a good system of axioms
for a notion of a ``tame subset of $\mathbb{R}^n$ '' seems to me to be the possibility
of proving such comparison theorems. So far I have
been content to draw out a plausible provisional system of axioms,
without any assurance that further axioms will not have to be
added, which probably only ``work on the actual material''\tnote{TN-034-01}{``Work on the actual material'' renders ``travail sur pièces'', preserving the author's expression for direct work on the material under study.}
will bring to light. The strongest of the axioms I have
introduced, and probably the one whose verification in specific
cases is (or will be) most delicate, is an \uline{axiom of}
\uline{triangulability} (tame, it goes without saying) of a tame subset of $\mathbb{R}^n$.
I have not attempted to prove the comparison theorem from these axioms
alone; nevertheless I have had the impression (rightly
or wrongly again !) that this proof, whether or not it requires
the introduction of some additional axiom, will present
no great technical difficulty. It is quite possible
that the technical difficulties in developing
satisfactory foundations for tame topology, including a
theory of dévissage of stratified tame structures, are
already concentrated for the most part in the axioms, and consequently
have essentially already been overcome by triangulability
theorems in the manner of Lojasiewicz and Hironaka. What is lacking,
once again, is by no means mathematicians' sometimes impressive
technical virtuosity, but the daring (or simply\frapparatus{V-FR-034-001}
the innocence ...) to free oneself from a familiar context
accepted by an unbroken consensus ...

The advantages of an axiomatic approach toward foundations
for tame topology seem fairly obvious to me. Thus,
to regard a complex algebraic variety, or the set of
real points of an algebraic variety defined over $\mathbb{R}$, as a
tame space, it seems preferable to work in the
``piecewise $\mathbb{R}$-algebraic'' theory, or even in the piecewise $\overline{\mathbb{Q}}_r$-algebraic
theory (where $\overline{\mathbb{Q}}_r = \overline{\mathbb{Q}} \cap \mathbb{R}$) when dealing with varieties defined

\margin{34}
\copytwo{A2-119-110-001}{\scancontext{https://grothendieck.umontpellier.fr/web/viewer.html?file=../119.pdf\#page=111}{\href{https://grothendieck.umontpellier.fr/web/viewer.html?file=../119.pdf\#page=111}{[copy 2, scan 119/110]} Copy 2 has handwritten “110” at bottom left.}}

\finishscan
% END PAGE esq_119_034
% BEGIN PAGE esq_119_035
\scanpage{119/35}
\mspage{33}

over number fields, etc. The introduction of a subfield
$K \subset \mathbb{R}$ associated with the theory $\hand{\mathcal{T}}$ (formed by the points of $\mathbb{R}$ that are
$T$-\uline{tame}, i.e. such that the corresponding one-point set
is tame) makes it possible to introduce, for every point $x$ of a tame space
$X$, a residue field $k(x)$, which is a subextension of $\mathbb{R}/K$
algebraically closed in $\mathbb{R}$, and of finite transcendence degree
over $K$ (bounded above by the topological dimension of $X$). When the
transcendence degree of $\mathbb{R}$ over $K$ is infinite, one obtains a
notion of the transcendence degree (or ``dimension'') of a point of a
tame space, close to the familiar notion in algebraic
geometry. Such notions are absent from ``semianalytic'' tame
topology, which, on the other hand, appears to be the most appropriate
topological context for including real
and complex analytic spaces.

Among the first theorems one expects in
a tame topology as I envisage it, apart from comparison
theorems\frapparatus{V-FR-035-001}, are statements establishing, in an
appropriate sense, the existence and uniqueness of ``the'' tubular neighborhood
of a closed tame subspace in a tame space (compact, for
simplicity), the concrete ways of obtaining it (for example from
any tame map\frapparatus{V-FR-035-002} $X \longrightarrow \mathbb{R}^{+}$ having $Y$ as
its zero set), the description of its ``boundary'' (although in
general it is by no means a manifold with boundary !) $\partial T$, which has
in $T$ a neighborhood isomorphic to the product of $T$ with an interval,
etc. Under suitable equisingularity hypotheses, one
expects $T$ to be endowed, in an essentially unique way,
with the structure of a locally trivial bundle over $Y$, having $\partial T$
as a subbundle. This is one of the least clear points in
my provisional intuition of the situation, whereas the
homotopy class of the anticipated structural map $T \longrightarrow Y$
has an obvious meaning, independently of any equisingularity
hypothesis, as a homotopy inverse of the inclusion map $Y \longrightarrow T$,
which must be a homotopy equivalence. One way of obtaining
such a structure a posteriori would be through the hypothetical equivalence of
isotopy categories envisaged at the beginning, taking into account the fact that

\margin{35}
\copytwo{A2-119-111-001}{\scancontext{https://grothendieck.umontpellier.fr/web/viewer.html?file=../119.pdf\#page=112}{\href{https://grothendieck.umontpellier.fr/web/viewer.html?file=../119.pdf\#page=112}{[copy 2, scan 119/111]} Copy 2 deletes the “a” in “résidual” and has a handwritten “e” in the left margin, giving “résiduel”.}}
\copytwo{A2-119-111-003}{\scancontext{https://grothendieck.umontpellier.fr/web/viewer.html?file=../119.pdf\#page=112}{\href{https://grothendieck.umontpellier.fr/web/viewer.html?file=../119.pdf\#page=112}{[copy 2, scan 119/111]} Copy 2 has a handwritten oblique stroke in the left margin beside the line containing “mais à part”.}}
\copytwo{A2-119-111-004}{\scancontext{https://grothendieck.umontpellier.fr/web/viewer.html?file=../119.pdf\#page=112}{\href{https://grothendieck.umontpellier.fr/web/viewer.html?file=../119.pdf\#page=112}{[copy 2, scan 119/111]} Copy 2 has handwritten “111” at bottom left.}}

\finishscan
% END PAGE esq_119_035
% BEGIN PAGE esq_119_036
\scanpage{119/36}
\mspage{34}

the functor $(U,Y,f) \longmapsto (X,Y)$ is defined in an obvious way,
independently of any theory of tubular neighborhoods\frapparatus{V-FR-036-002}.

It will probably be said, not without some reason, that all
this is perhaps only dreams, which will vanish into smoke as soon
as one attempts detailed work, or even before that,
in the face of certain known or quite obvious facts that may have
escaped me. Certainly, only work on the actual material will make it possible to
separate the true from the false and to know the actual substance. The
only thing in all this about which I have no
doubt is the need for such work on foundations, in
other words, the artificial nature of the present foundations of
topology, and of the difficulties they raise at every
step. On the other hand, it is quite possible that my formulation
of a theory of dévissage of stratified structures, as
an equivalence theorem for suitable isotopy (or even
$\infty$-isotopy) categories, is too optimistic. I should add,
however, that I also have scarcely any doubt that the theory
of these dévissages that I developed two years ago, although
it remains partly heuristic, does indeed express a reality
that is as tangible as can be. In part of my work,
for want of a ready-made ``tame'' context, and
in order nonetheless to have precise and provable statements,
I was led to postulate on the initial stratified structure
some entirely plausible additional structures,
in the nature of given local retractions in particular,
which then do indeed allow the construction of a canonical system
of spaces, parametrized by the ordered set of ``flags'' Drap(I)
of the ordered set $I$ indexing the strata, these spaces playing
the role of the spaces $(U,Y)$ from earlier, linked to one another by
embedding maps and proper fibrations, which make it
possible to reconstruct, in just as canonical a way\frapparatus{V-FR-036-001}, the initial
stratified structure, including those ``additional structures''\notecall{copy1-7}{call-esq_119_036-7-1}{(${}^{7}$)}.
The only trouble is that these latter seem to be a superfluous element of
structure, which is by no means given in ordinary geometric

\margin{36}
\copytwo{A2-119-112-001}{\scancontext{https://grothendieck.umontpellier.fr/web/viewer.html?file=../119.pdf\#page=113}{\href{https://grothendieck.umontpellier.fr/web/viewer.html?file=../119.pdf\#page=113}{[copy 2, scan 119/112]} Copy 2 has handwritten “112” at bottom left.}}

\finishscan
% END PAGE esq_119_036
% BEGIN PAGE esq_119_037
\scanpage{119/37}
\mspage{35}

situations, for example for the compact moduli space
$M^{\wedge}_{g,\nu}$ with its canonical ``stratification at infinity'', given
by the Mumford-Deligne normal crossings divisor. Another
difficulty, probably less serious, is that the so-called
moduli ``space'' is in fact a \uline{multiplicity} --- technically,
this is expressed above all by the need to replace
the index set $I$ for the strata by an (essentially finite)
index \uline{category}, in this case that of
``MD graphs'', which ``parametrize'' the possible ``combinatorial structures''
of a stable curve of type $(g,\nu)$. That said, I can
affirm that the general theory of dévissage, developed especially
under the pressure of the needs of \uline{this} cause, has indeed proved
a valuable guide, leading to a progressive understanding,
with unfailing coherence, of certain essential aspects
of the Teichmüller tower (that is, essentially
of the ``structure at infinity'' of the ordinary Teichmüller groups).
It is this approach that finally led me, in the following
months, toward the principle of a purely combinatorial
construction of the tower of Teichmüller groupoids, in the spirit
sketched \unsure{.}above\tnote{TN-037-01}{The scan has a doubtful dot between ``plus'' and ``haut'' in ``plus haut''; it is retained beside the corresponding English word.} (cf. par. 2).

Another satisfactory test of coherence comes from the
``topossic'' viewpoint.\tnote{TN-037-02}{``Topossic'' renders the author's adjective ``topossique'', referring here to the context of topoi.} Indeed, since my interest in moduli
multiplicities stems above all from their algebraic-geometric
and arithmetic meaning, it is the \uline{algebraic} moduli
multiplicities, over the absolute base field $\mathbb{Q}$, that have
primarily interested me, and a dévissage ``at infinity'' of their
geometric fundamental groups (i.e. the \uline{profinite} Teichmüller
groups) compatible with the natural actions of
$\mathrm{I}\!\Gamma = \mathrm{Gal}(\bar{\mathbb{Q}}/\mathbb{Q})$. This seemed from the outset to exclude the possibility of
referring to a hypothetical theory of dévissage of stratified
structures in a context of ``tame topology'' (or even of
ordinary topology, after a fashion), except as an
entirely heuristic guiding thread. The question then arose of translating,
into the context of the topoi (in this case the étale topoi)

\margin{37}

\finishscan
% END PAGE esq_119_037
% BEGIN PAGE esq_119_038
\scanpage{119/38}
\mspage{36}

involved in the situation, the theory of dévissage that
I had reached\frapparatus{V-FR-038-001} in an entirely different context --- with the
additional task, subsequently, of drawing out a general comparison
theorem, on the model of the well-known theorems, to
compare the invariants obtained (in particular the homotopy types of
various tubular neighborhoods) in the transcendental setting and in
the scheme-theoretic setting. I was able to convince myself that such a formalism
of dévissage did indeed make sense in the (so-called
``abstract'' !) context of general topoi, or at least of noetherian
topoi (such as those that arise here), through a suitable
notion of the \uline{canonical tubular neighborhood of a subtopos} in an
ambient topos. Once this notion has been acquired, with certain simple
formal properties, the description of the ``dévissage'' of a
stratified topos is considerably simpler even in this
setting than in the topological (tame) setting. It is true that
here too work on foundations remains to be done, in particular on
the very notion of the tubular neighborhood of a subtopos --- and it is
surprising, moreover, that this work (as far as I know)
has still not been done, that is, that no one (in the more
than twenty years that a context of étale topologies has existed) seems
to have needed it; surely a sign that the understanding
of the topological structure of schemes has not progressed all that
much since the work of Artin-Mazur ...

Once the double task of preliminary work
(more or less heuristic) around the notion of the dévissage of a
stratified space or topos had been accomplished, which was a crucial stage in
my understanding of moduli multiplicities, it moreover
became apparent that for the needs of the latter, one can probably
short-circuit at least a good part of this theory by
direct geometric arguments. The fact remains that for
me, the formalism of dévissage that I reached has proved
its usefulness and coherence, independently of any
question about the most appropriate foundations that would allow
it to be given its full meaning.

\margin{38}

\finishscan
% END PAGE esq_119_038
% BEGIN PAGE esq_119_039
\scanpage{119/39}
\mspage{37}

6. One of the most interesting foundational theorems of (tame)
topology that ought to be developed would be a theorem
of ``dévissage'' (again !) of a proper tame map
of tame spaces,\frapparatus{V-FR-039-001}

\[ f : X \longrightarrow Y , \]

through a decreasing filtration of $Y$ by closed tame
subspaces $Y^{i}$, such that over the ``open strata''
$Y^{i} \setminus Y^{i-1}$ of this filtration, $f$ induces a locally
trivial fibration (from the tame viewpoint, it goes without saying). Such
a statement should further be generalized and made more precise in various
ways, in particular by requiring the existence of an analogous
\uline{simultaneous} dévissage for $X$ and a given finite family of closed
(tame) subspaces of $X$. The very notion of a locally
trivial fibration in the tame sense can also be strengthened
considerably, taking into account the fact that the open
strata $U_{i}$ are \uline{better} than spaces with a purely local tame
structure, because they are obtained as the difference of
two tame spaces, compact when $Y$ was compact. Between the
notion of a compact tame space (which can be realized as one of the
initial ``models'' in an $\mathbb{R}^{n}$) and that of a ``locally
tame'' (locally compact) space, which follows from it in a fairly obvious
way, there is a somewhat more delicate notion of a \uline{``globally}
\uline{tame'' space} $X$, obtained as the difference $X^{\wedge}\setminus Y$ of two compact
tame spaces, it being understood that no distinction is made between
the space defined by a pair $(X^{\wedge},Y)$ and that defined by a
pair $(X'^{\wedge},Y')$ obtained from it by a tame
(necessarily proper) map

\[ g : X'^{\wedge} \longrightarrow X^{\wedge} \]

inducing a bijection $g^{-1}(X) \xrightarrow{\ \sim\ } X$, taking $Y' = g^{-1}(Y)$.
Perhaps the most interesting natural example is that in which one
starts from a separated scheme of finite type over $\mathbb{C}$ or over $\mathbb{R}$, taking

\margin{39}

\finishscan
% END PAGE esq_119_039
% BEGIN PAGE esq_119_040
\scanpage{119/40}
\mspage{38}

for $X$ its set of complex or real points, which inherits
a global tame structure by means of the scheme-theoretic
compactifications\frapparatus{V-FR-040-001} (which exist by Nagata) of the initial scheme.
This notion of a globally tame space is associated with a notion
of a \uline{globally tame map}, which in turn makes it possible to
strengthen accordingly the notion of a locally trivial
fibration, in the statement of a dévissage theorem for a
map $f : X \longrightarrow Y$ (no longer necessarily proper) in
the context of globally tame spaces.

Last summer Zoghman Mebkhout informed me
that a dévissage theorem in this spirit had recently been obtained
in the context of real and/or complex analytic spaces,
with $Y^{i}$ that, this time, are analytic subspaces
of $Y$. This result makes it plausible that we already have
sufficiently powerful technical means to prove
a dévissage theorem in the tame context as well, apparently
more general, but probably less arduous.

It is also the context of a tame topology
that should, it seems to me, make it possible to formulate precisely
a very reliable general principle that I have long used in
a great many geometric situations, which I call the
``\uline{principle of harmless choices}'' --- as useful as it appears vague !
It says that when, for the needs of some construction
of one geometric object in terms of others, one is led to make
a certain number of arbitrary choices along the way, so
that the object obtained apparently depends on these choices, and is
therefore marred by a lack of canonicity, this defect is indeed
serious (and removing it requires a more careful analysis of
the situation, of the notions used, of the data introduced, etc.)
whenever at least one of these choices is made in a
``space'' that is not ``contractible'', i.e. whose $\pi_{0}$ \uline{or} one of
the higher invariants $\pi_{i}$ is nontrivial; whereas this defect is
merely apparent, the construction is ``essentially
canonical'' and will not really cause trouble, whenever

\margin{40}

\finishscan
% END PAGE esq_119_040
% BEGIN PAGE esq_119_041
\scanpage{119/41}
\mspage{39}

the choices made are all ``harmless'', i.e. are made
in \uline{contractible} spaces. When one tries, in particular
cases, to pin down this principle more closely, it seems that one comes
each time to the notion of ``$\infty$-isotopy categories'' expressing
a given situation, finer than the more naïve isotopy
(= 0-isotopy) categories, obtained by retaining only the $\pi_0$
of the spaces of isomorphisms that arise in the situation,
whereas the $\infty$-isotopy viewpoint retains their entire
homotopy type. For example, the naïve isotopy viewpoint for
oriented compact surfaces with boundary of type $(g,\nu)$ is ``good''
(with no hidden boomerang !) exactly in the cases that I call
``anabelian'' (and that Thurston calls ``hyperbolic''), i.e. distinct
from $(0,0)$, $(0,1)$, $(0,2)$, $(1,0)$ --- which are also precisely the cases
in which the automorphism group of the surface has a
\uline{contractible} identity component. In the other cases, except
the case $(0,0)$ of the sphere with no holes, it suffices to work with
1-isotopy categories to express the essential geometric-topological
facts satisfactorily by algebraic means, since
the said connected component is then a $K(\pi,1)$. Working
in a 1-isotopy category moreover amounts to working
in a bicategory, i.e. with $\underline{\mathrm{Hom}}(X,Y)$ that are (no longer
discrete sets as in the 0-isotopy viewpoint,
but) groupoids (whose $\pi_0$ are precisely the
0-isotopy Hom). It is the description in purely algebraic
terms of this bicategory that is given in the final part
of Yves Ladegaillerie's thesis (cf. par. 3).

If I have dwelt here at greater length on the theme
of the foundations of tame topology, which is by no means one of
those to which I intend to devote myself as a priority in the
coming years, it is probably precisely because I feel that
there is all the more a cause here that needs to be pleaded, or
rather: work of great current relevance that needs hands !
As formerly with new foundations for algebraic
geometry, it is not pleas that overcome the inertia

\margin{41}

\finishscan
% END PAGE esq_119_041
% BEGIN PAGE esq_119_042
\scanpage{119/42}
\mspage{40}

of acquired habits, but tenacious, meticulous work,
probably over a long period, yielding eloquent harvests
day by day.

I would like to say a few more words about an earlier
reflection (late 60s ?), very close to the one just
discussed, inspired by Nash's ideas, which had
greatly struck me. Here, instead of defining axiomatically a notion
of a ``tame theory'' by giving a notion of a ``tame subset
of $\mathbb{R}^n$ '' satisfying certain conditions (chiefly of stability),
what I was after was an axiomatization of the notion of a ``smooth manifold'' and of the
differentiable formalism on such manifolds,
by giving, for each natural number $n$, a subring
$\unsure{\sourceglyph{119A}}_n$ of the ring of germs of real functions at the origin in $\mathbb{R}^n$.\tnote{TN-042-01}{The visible typed symbol used for this system of rings is reproduced from the scan. Its alphabet remains unresolved; the doubt marks preserve that uncertainty without assigning a letter family.}
These are the functions that will be admitted for expressing the ``changes
of chart'' for the corresponding notion of a $\unsure{\sourceglyph{119A}}$-manifold, and the
first task was to draw out\frapparatus{V-FR-042-001} a system of axioms on the
system $\unsure{\sourceglyph{119A}} = (\unsure{\sourceglyph{119A}}_n)_{n\in\underline{N}}$ that ensures this notion of a manifold a
flexibility comparable to that of a $C^\infty$, or real analytic
(or Nash), manifold. Depending on the kind of familiar constructions that
one wants to be able to carry out in the context of $\unsure{\sourceglyph{119A}}$-manifolds, the
relevant system of axioms is more or less restricted, or rich.
Very little suffices when it is only a matter of developing the
differential formalism, with the construction of jet bundles, De Rham
complexes, etc. To obtain a statement of the type ``quasi-finite implies
finite'' (for a map in a neighborhood of a point), which
emerged as a key statement in the local theory of analytic
spaces, one needs a stability axiom of a more delicate nature,
in Weierstrass's ``Vorbereitungssatz''\tnote{TN-042-02}{The German term ``Vorbereitungssatz'' means ``preparation theorem''.} (*). In other

\rule{4cm}{0.4pt}

(*) It may seem simpler to say that the (local) rings
$\unsure{\sourceglyph{119A}}_n$ are \uline{henselian}, which is equivalent. But in this latter
form it is by no means clear a priori that the condition in
question is in the nature of a stability condition, an
important circumstance, as will emerge in the reflections that follow.

\margin{42}

\finishscan
% END PAGE esq_119_042
% BEGIN PAGE esq_119_043
\scanpage{119/43}
\mspage{41}

questions, an axiom of stability under analytic continuation
(in $\mathbb{C}^{n}$) appears necessary. The most draconian axiom that
I have been led to introduce, also a stability axiom,
concerns the integration of Pfaff systems, ensuring that certain
Lie groups, or even all of them, are $\unsure{\sourceglyph{119A}}$-manifolds. In all
this, \uline{I have taken care not to assume that the} $\unsure{\sourceglyph{119A}}_{n}$ \uline{are}
\uline{$\mathbb{R}$-algebras}, so a constant function on an $\unsure{\sourceglyph{119A}}$-manifold is
``admissible'' only when its value belongs to a certain subfield
$K$ of $\mathbb{R}$ (this is, if one wishes, $\unsure{\sourceglyph{119A}}_{0}$). This subfield may very well
be $\mathbb{Q}$, or its algebraic closure $\bar{\mathbb{Q}}_{r}$ in $\mathbb{R}$, or any other
subextension of $\mathbb{R}/\mathbb{Q}$, preferably even of finite,
or at least countable, transcendence degree over $\mathbb{Q}$. This makes it possible,
for example, as earlier for tame spaces, to associate
with every point $x$ of a manifold (of type $\unsure{\sourceglyph{119A}}$) a residue
field $k(x)$, which is a subextension of $\mathbb{R}/K$. A fact that
seems important to me here is that, even in its strongest form,
the system of axioms does \uline{not} imply that one must have $K = \mathbb{R}$. More
precisely, since \uline{all} the axioms are stability
axioms, it follows that for a given set $S$ of germs
of real analytic functions at the origin (in various spaces\frapparatus{V-FR-043-001}
$\mathbb{R}^{n}$), there exists a smallest theory $\unsure{\sourceglyph{119A}}$ for which these germs
are admissible, and that this theory is ``countable'', i.e. the $\unsure{\sourceglyph{119A}}_{n}$
are countable, as soon as $S$ is countable. A fortiori, $K$ is then
countable, i.e. of countable transcendence degree over $\mathbb{Q}$.

The idea here is to introduce, by means of this
axiomatic approach, a notion of an ``\uline{elementary}'' (real analytic) function,
or rather, a whole hierarchy of such notions. For a function
of 0 variables, i.e. a constant, this notion gives that
of an ``elementary constant'', including in particular (in the case of
the strongest system of axioms) constants such as $\pi$, $e$, and
a multitude of others, by taking values of admissible functions
(such as the exponential, the logarithm, etc.) for systems
of ``admissible'' values of the argument. One feels that the relationship
between the system $\unsure{\sourceglyph{119A}} = (\unsure{\sourceglyph{119A}}_{n})_{n\in\underline{N}}$ and the field of rationality $K$

\margin{43}

\finishscan
% END PAGE esq_119_043
% BEGIN PAGE esq_119_044
\scanpage{119/44}
\mspage{42}

corresponding to it must be very close, at least for $\unsure{\sourceglyph{119A}}$ that
can be generated by a finite ``system of generators'' S ---
but it is to be feared that the slightest interesting question one
might ask about this situation is currently out of
reach \notecall{copy1-1}{call-esq_119_044-1-2}{($^1$)}.

These earlier reflections regained some relevance
for me with my later reflection on tame theories.
It seems to me, indeed, that it is possible to associate naturally
with a ``differentiable theory'' $\unsure{\sourceglyph{119A}}$ a tame theory $\hand{\mathcal{T}}$
(probably having the same field of constants), in such a way that
every $\unsure{\sourceglyph{119A}}$-manifold\frapparatus{V-FR-044-001} is automatically endowed with a $\hand{\mathcal{T}}$-tame structure,
and conversely that for every compact $\hand{\mathcal{T}}$-tame space $X$, one can
find a closed nowhere dense tame subset $Y$ in $X$, such that $X \smallsetminus Y$
comes from an $\unsure{\sourceglyph{119A}}$-manifold, and that moreover this $\unsure{\sourceglyph{119A}}$-manifold
structure is unique at least in the following sense: two such
structures coincide in the complement of a nowhere dense tame
subset $Y' \supset Y$ of $X$. The theory of dévissage of
stratified tame structures (discussed in the preceding paragraph),
in the case of smooth strata, should moreover raise
much more precise questions still about the comparison of tame
structures with structures of differentiable type (or rather,
$\mathbb{R}$-analytic type). I suspect that the kind of axiomatization proposed
here for the notion of a ``differentiable theory'' would provide a
natural setting for formulating such questions with all the precision
and generality one could wish for.


7. Since March of last year, thus for nearly
a year, the greater part of my energy has been devoted to
reflection on the \uline{foundations of noncommutative} (\uline{co})\uline{homological}
\uline{algebra}, or what ultimately amounts to the same thing, \uline{homotopical}
\uline{algebra}. These reflections have taken concrete form in a voluminous
bundle of typewritten notes, intended to form the first
volume (currently nearing completion) of a work in two

\margin{44}

\finishscan
% END PAGE esq_119_044
% BEGIN PAGE esq_119_045
\scanpage{119/45}
\mspage{43}

volumes to be published by Hermann, under the common title ``\uline{Pursuing}
\uline{Stacks}''. I currently expect (after successive expansions
of the initial undertaking) that the manuscript\frapparatus{V-FR-045-001} of the two
volumes together, which I hope to finish during the year so as not
to have to return to it, will run to about 1,500 typewritten pages.
These two volumes are moreover, for me, the first in a
larger series, under the common title ``\uline{Mathematical Reflections}'',
in which I intend to develop to some extent certain of the themes sketched
in the present report.

Since this is a work currently being written,
and indeed approaching completion, whose first volume will probably appear
this year and will contain a detailed introduction, it
is probably less interesting for me to dwell here on this theme
of reflection, and I shall therefore be content to speak of it very briefly.
This work seems to me somewhat marginal in relation to the themes
I have just sketched, and does not represent (it seems to me)
a real renewal of outlook or approach in relation to
my mathematical interests and vision before 1970. If I
suddenly resolved to undertake it, this was almost out of desperation,
when nearly twenty years have elapsed since a number of visibly
fundamental questions, ripe to be carried through to completion,
were posed in quite clear terms, without
anyone seeing them or taking the trouble to probe them. Even
today, the basic structures involved in the
homotopy viewpoint in topology, including even commutative
homological algebra, are not understood, and to my
knowledge, after the work of Verdier, Giraud, and Illusie on
this theme (which constitute so many ``kick-offs'' still
awaiting a continuation ...), there has been no effort in this direction. I
should probably make an exception for\unsure{.}\tnote{TN-045-01}{A visible baseline dot follows ``pour'' in the scan. Its punctuation status remains uncertain; the mark is retained after ``for'' here.} the axiomatization
carried out by Quillen on the notion of a model category, in
the late 60s, and taken up in various forms by
various authors. This work, at the time and still today, has
greatly attracted me and taught me much, while going in a rather

\margin{45}

\finishscan
% END PAGE esq_119_045
% BEGIN PAGE esq_119_046
\scanpage{119/46}
\mspage{44}

different direction from the one that was and is close to my heart. It certainly
introduces derived categories in various noncommutative
contexts, but without entering into the question of the essential internal
structures of such a category, also left open
in the commutative case by Verdier and, after him, by Illusie. Likewise,
the question of identifying the natural ``coefficients''
for a noncommutative cohomological formalism, beyond
the stacks (which should be called 1-stacks) studied in Giraud's
book, remained open --- or rather, the rich and
precise intuitions that answer it, drawn from many examples
coming in particular from algebraic geometry, still await
a precise and flexible language to give them form.

I return to certain aspects of these foundational
questions in 1975, on the occasion (if I remember correctly) of a
correspondence with Larry Breen (two letters from this correspondence
will be reproduced as an appendix to Chap. I of volume 1,
``History of Models'',\tnote{TN-046-01}{The French volume title is ``Histoire de Modèles''. ``Histoire'' can mean either a history or a story; ``History of Models'' preserves the singular title.} of Pursuing Stacks). At this point
the intuition emerges that $\infty$-groupoids must constitute
particularly adequate models for homotopy types,
n-groupoids corresponding to \uline{truncated} homotopy types
(with $\pi_i = 0$ for $i > n$). This same intuition, by very
different paths, was rediscovered by Ronnie Brown in Bangor and
some of his students, but using a rather restrictive notion of an $\infty$-groupoid
(which, among 1-connected homotopy types,
models only products of Eilenberg-Maclane spaces\frapparatus{V-FR-046-001}). It was
stimulated by a discursive correspondence with Ronnie Brown
that I finally resumed reflection, beginning with an attempt
to define a broader notion of an $\infty$-groupoid (subsequently renamed
``stack in $\infty$-groupoids'' or simply ``stack'', with the
understanding: over the punctual topos), which, one thing leading to another,
brought me to Pursuing Stacks. The volume ``History of Models''
moreover constitutes an entirely unforeseen digression
from the initial undertaking (the famous stacks being temporarily
forgotten, and not expected to reappear until around
page 1,000 ...).

\margin{46}

\finishscan
% END PAGE esq_119_046
% BEGIN PAGE esq_119_047
\scanpage{119/47}
\mspage{45}

This work is not entirely isolated in relation to
my more recent interests. For example, my reflection on the
moduli multiplicities $M^{\wedge}_{g,\nu}$ and their stratified structure
revived reflection on a van Kampen theorem of dimension
$>$ 1 (also one of the favorite themes of the Bangor
group), and may have helped prepare the ground for
the work of greater scope the following year. The latter
also connects at times with a reflection dating from the same
year 1975 (or the following year) on a ``De Rham complex with
divided powers'', which was the subject of my last public
lecture, at IHES\tnote{TN-047-01}{IHES is the Institut des Hautes Études Scientifiques; see the \href{https://www.ihes.fr/en/}{institute's website}.} in 1976, and whose manuscript, entrusted
to someone I no longer remember after the lecture, is moreover lost.
It was at the time of this reflection that the intuition
of a ``schematization'' of homotopy types also began to germinate, which,
seven years later, I am trying to make precise in a (particularly
hypothetical) chapter of History of Models.

The work of reflection undertaken in Pursuing
Stacks is a little like a debt I would be discharging
toward a scientific past in which, for some fifteen
years (between 1955 and 1970), the development of cohomological
tools was the constant leitmotif in my work on the
foundations of algebraic geometry. If the present resumption of this
theme has taken on unexpected dimensions, it is not,
however, out of pity for a past, but because of the many unforeseen
things constantly breaking in, mercilessly upsetting the
plans and intentions envisaged --- a little like a tale from the Thousand and
One Nights, in which one's attention is kept in suspense through
twenty other tales before learning the final word of the first.


8. I have as yet said very little about the more down-to-earth
reflections on two-dimensional topological geometry, associated
in particular with my activities as a teacher and those called ``research
supervision''. On several occasions, I have seen\unsure{.}\tnote{TN-047-02}{The doubtful dot after ``vu'' is retained after the corresponding ``seen''.} opening up before

\margin{47}

\finishscan
% END PAGE esq_119_047
% BEGIN PAGE esq_119_048
\scanpage{119/48}
\mspage{46}

me vast and rich fields ripe for harvest, without
ever managing to communicate this vision, and the spark
that accompanies it, to a male (or female) student of mine, or to make it
lead to a common exploration of greater or lesser
duration. Each time, right up to the present day, after an exploration
of a few days or a few weeks, in which I discovered,
as a scout, riches unsuspected at the outset, the journey
was cut short when it became clear that I would be alone in
continuing it. Stronger interests then took precedence over
a journey that consequently appeared as a digression, or even
a scattering of effort, rather than an adventure pursued together.

One of these themes was that of plane polygons, centered
on the moduli varieties that can be associated with them. One of
the surprises here was the irruption of algebraic geometry into
a context that had seemed to me quite far removed from it. This kind of
surprise, linked to the ubiquity of geometry in geometry as
such, has moreover recurred on several occasions.\tnote{TN-048-01}{The source repeats ``géométrie'' without an added adjective and has the form ``répéré''. ``Recurred'' follows the repetition expressed by ``à plusieurs reprises''; the intended French verb remains uncertain, and the checked French wording is retained.}

Another theme was that of curves (in particular
circles) immersed in a surface, with particular attention
to the ``stable'' case in which the singular points are ordinary
double points (and also the more general case in which the
different branches at a point cross one another), often with
the additional assumption that the immersion is ``cellular'',
i.e. gives rise to a map. A variant of situations of
this type is that of immersions of a surface with nonempty boundary,
and first and foremost of a disk (which had been pointed out to me by
A. Campo some ten years ago). Beyond the question of
various combinatorial formulations of such situations, which
is scarcely more than an exercise in syntax, I have been interested
above all in a dynamic view of the possible configurations,
with\unsure{.}\tnote{TN-048-02}{The doubtful dot after ``avec'' is retained after ``with''.} the passage from one to another by continuous deformations, which
can be decomposed into composites of two types of \uline{elementary}
\uline{operations} and their inverses, namely the ``\uline{sweeping}'' of a branch
of a curve across a double point, and the \uline{erasure} or \uline{creation}

\margin{48}

\finishscan
% END PAGE esq_119_048
% BEGIN PAGE esq_119_049
\scanpage{119/49}
\mspage{47}

\uline{of a bigon}. (The first of these operations also plays a
key role in a "dynamic" theory of systems of pseudolines
in a real projective plane.) One of the first questions
that arise here is to determine the different \uline{immersion
classes} of a circle or a disk (say) modulo these elementary
operations; another is to see which immersions
of the boundary of the disk come from an immersion of the
disk, and to what extent the former determine the latter.
Here again, it seemed to me that a systematic study
of the relevant moduli varieties (of infinite dimension in this
case, unless one manages to give a purely combinatorial
version of them) should provide the most effective ``focus'',
in a sense forcing us to ask ourselves the most pertinent
questions. Unfortunately, reflection on even the
most obvious and down-to-earth questions has remained embryonic.
As the only tangible result, I can point to a
theory of canonical ``dévissage'' of a stable cellular immersion
of the circle in a surface into ``indecomposable'' immersions,
by ``telescoping'' such immersions. Unfortunately, I have not
succeeded in seeing my insights into the question turn into
a DEA research project, nor other insights (into a complete
theoretical description, in terms of fundamental groups of
topological 1-complexes, of the immersions of a surface with boundary
that extend a given immersion of its boundary) into the beginning of a
state doctorate thesis\tnote{TN-049-01}{``State doctorate'' renders ``doctorat d'état''. The later ``third-cycle theses'' renders ``thèses de troisième cycle''; the source uses these degree designations.}...

A third theme, pursued simultaneously for three
years at various teaching levels (from the elective for
first-year students to three third-cycle theses currently
being pursued on this theme), concerns the topological-combinatorial
classification of systems of lines or pseudolines.
On the whole, my students' participation here has been less
disappointing than elsewhere, and I have sometimes had the pleasure of learning
from them interesting things that would not have occurred to me.
Our joint reflection has of necessity, however, been limited

\margin{49}

\finishscan
% END PAGE esq_119_049
% BEGIN PAGE esq_119_050
\scanpage{119/50}
\mspage{48}

to a very elementary level. Recently I finally
devoted a month of intensive reflection to developing a
purely combinatorial construction of a kind of ``moduli surface''
associated with a system of $n$ pseudolines, which classifies the
different possible ``relative positions'' (stable or otherwise) of an
$(n+1)$th pseudoline with respect to the given system, or again:
the different possible ``affinizations'' of this system, through the
different possible choices of a ``pseudoline at infinity''. I
have the impression of having put my finger on a remarkable object,
bringing to light an unforeseen order in classification
questions that until now appeared rather chaotic !
But the present report is not the place to dwell further
on this subject.

Since 1977, in all questions (as in the
last two themes I have just mentioned) involving
two-dimensional maps, the possibility of realizing them canonically
on a conformal surface, hence on a complex algebraic
curve in the oriented compact case, remains a constant watermark
in my reflection. In practically all cases (in fact, all
cases except that of certain spherical maps with ``few automorphisms''),
such a conformal realization actually implies a
\uline{canonical riemannian metric}, or at least one canonical up to a
multiplicative constant. These new elements of structure
(even without taking into account the arithmetic element discussed
in par. 3) are such as to transform profoundly
the initial appearance of the questions addressed, and the methods of
approach. The beginnings of familiarity with Thurston's beautiful ideas
on the construction of Teichmüller space, in terms of a very
simple game of hyperbolic riemannian surgery, confirm me
in this presentiment. Unfortunately, the very modest level of knowledge
of almost all the students who have worked with me during\frapparatus{V-FR-050-001}
these last ten\ \ years does not allow me to broach such possibilities
with them, even by allusion, whereas
the assimilation of a minimal combinatorial language already encounters,

\margin{50}

\finishscan
% END PAGE esq_119_050
% BEGIN PAGE esq_119_051
\scanpage{119/51}
\mspage{49}

quite often, considerable psychological obstacles. That is why,
in certain respects and increasingly in recent years,
my teaching activities have often acted as a burden,
rather than as a stimulus to the development of geometric
reflection that is at all advanced, or merely delicate.

9. This seems an opportune occasion to make a brief
appraisal of my teaching activity since 1970, that is, since
it has taken place in a university setting. This contact
with a very different reality has taught me a great
deal, of an order quite different, moreover, from the
merely pedagogical or scientific. This is not the place to
dwell on that subject. I also spoke at the beginning of this report of the
role played by this change of professional milieu in the
renewal of my approach to mathematics, and of my
mathematical centers of interest. On the other hand, when I take stock
of my teaching activity at the level of research properly
speaking, I arrive at a finding of clear and unambiguous failure. In the more
than ten years that this activity has continued year after year within
the same university institution, I have not, at any time, managed
to bring about a place where ``something is happening'' --- where something
``passes'',\tnote{TN-051-01}{The French repeats the verb ``passer'': ``il se passe quelque chose'' means that something is happening, and ``quelque chose passe'' says that something passes among the people. The second literal phrasing retains that verbal echo.} among a group, however small, of people
linked by a common adventure. On two occasions, it is true, around
the years 74 to 76, I had the pleasure and privilege of inspiring
a student to undertake work of scope, pursued with enthusiasm:
Yves Ladegaillerie's work mentioned earlier (par. 3) on
questions of isotopy in dimension 2, and Carlos Contou-Carrère's
(whose mathematical passion had not waited for his encounter
with me to blossom) unpublished work on local and global
Jacobians over general base schemes (part of which
was announced in a note in CR).\tnote{TN-051-02}{CR abbreviates \emph{Comptes Rendus} of the Académie des sciences; see the \href{https://comptes-rendus.academie-sciences.fr/page/history-comptes-rendus/}{journal's historical account}.} Apart from these two cases,
my role over these ten years has been confined to passing on,
as best I could, the rudiments of the mathematician's craft to some
twenty students\unsure{.}\tnote{TN-051-03}{The doubtful dot after ``élèves'' is retained after ``students''.} at the research level, or at least to

\margin{51}

\finishscan
% END PAGE esq_119_051
% BEGIN PAGE esq_119_052
\scanpage{119/52}
\mspage{50}

those among them who persevered long enough with me, reputed
to be more demanding than others, to produce a first acceptable
piece of work in black and white (sometimes also work better
than acceptable and more than a single piece of work\frapparatus{V-FR-052-001}, done with taste\ \ and
carried through to the end). Given the circumstances, even among the few who
persevered, fewer still will have the opportunity
to practice this craft and thereby, while earning their bread,
to deepen it.

10. Since last year I have felt that, in the course of my
activity as a university teacher, I have learned all I had
to learn from it and taught all I can teach there, and
that it has ceased to be truly useful, to me as well as to
others. Persisting in continuing it under these conditions would
seem to me a waste, both of human resources and of public
funds. That is why I have requested secondment to the CNRS (which
I had left in 1959 as a newly minted research director,
to enter IHES). I am aware, moreover, that the employment
situation is tight at the CNRS as elsewhere, that the outcome of my
application is doubtful, and that a position awarded to me would
be at the expense of a younger researcher who would remain without one.
But it is also true that this would free my position at USTL\tnote{TN-052-01}{USTL is the Université des Sciences et Techniques du Languedoc; the expansion is given in the \href{https://grothendieck.umontpellier.fr/}{archive's list of abbreviations}.}
for someone else. That is why I have no scruples about
making this application, and, should it be necessary, about returning to the charge
if it is not accepted this year.

In any event, this application will have given me
the occasion to write this sketch of a program, which otherwise would
probably never have seen the light of day. I have tried to be brief without being
sibylline and also, afterward, to make it easier and
more attractive to read by adding a table of contents. If despite
this it may seem long for the occasion, I apologize.
It seems short to me for its contents, knowing well that ten years of
work would not be too much to follow through to the end even the least

\margin{52}

\finishscan
% END PAGE esq_119_052
% BEGIN PAGE esq_119_053
\scanpage{119/53}
\mspage{51}

of the themes sketched (assuming there is an ``end''...), and
a hundred years would be little for the richest of them !

Behind the apparent disparity of the themes mentioned here,
an attentive reader will perceive, as I do, a deep unity. This
manifests itself in particular through a common source of inspiration,
the geometry of surfaces, present in all these themes, most
often in the foreground. This source, in relation to my mathematical
``past'', represents a renewal, but by no means a
break. Rather, it points the way toward a new approach
to that still mysterious reality, that of ``\uline{motives}'', which
fascinated me more than any other in the last years of that
past (* ). This fascination has by no means vanished;
rather, it forms part of that of the most burning for me of all
the themes mentioned earlier. But today I am no longer,
as I once was, the willing prisoner of endless tasks,
which so often prevented me from launching into the unknown,
mathematical or otherwise. The time of \uline{tasks} is over for me. If
age has brought me anything, it is to be lighter.

January 1984

(* ) On this subject, see my comments in the ``Thematic Sketch''
of 1972 attached to the present report, under the final heading
``motivic ramblings'',\tnote{TN-053-01}{The French titles are ``Esquisse Thématique'' and ``divagations motiviques''.} (loc. cit. pages 17--18).

\margin{53}

\finishscan
% END PAGE esq_119_053
% BEGIN PAGE esq_119_054
\scanpage{119/54}
\mspage{52}

\notedef{copy1-1}{(\textsuperscript{1})} The expression ``out of reach'' here (and again further on for
an entirely different question), which I let pass despite
a reluctance, now strikes me as decidedly hasty and
unfounded. I have already been able to observe on other occasions that when
oracles (here myself !) declare with a knowing
(or doubtful) air that such-and-such a problem is ``out of reach'', this is
at bottom an entirely subjective assertion. It means
simply, apart from the fact that the problem is supposed not to have been
solved yet, that the speaker is short of ideas on the
question, or, probably more precisely, that he faces
it without feeling or enthusiasm, that it ``does nothing for him'', and
that he has no desire to do anything with it --- which
is often reason enough to want to discourage others
from doing so. This did not prevent the beautiful and lamented conjectures
of Mordell, Tate, and Chafarévitch, like Mr. de la Palisse,\tnote{TN-054-01}{The allusion evokes the proverbially obvious observation of a man still alive moments before dying. Here the joke concerns conjectures still declared out of reach just before they succumb. On the French allusion, see the entry ``truisme'' in the \href{https://oqlf.gouv.qc.ca/ressources/bibliotheque/dictionnaires/VocabulairesPDF/terminologie_figuresdestyle.pdf}{OQLF vocabulary of figures of speech}.}
from still being reputed ``out of reach'' at the very moment
of succumbing, the poor things ! --- Moreover, already in the days
following the writing of the present report, which brought me back
into contact with questions from which I had somewhat distanced myself
over the past year, I became aware of a new remarkable
property of the outer action of an absolute Galois
group on the fundamental group of an algebraic curve, which had
escaped me until now\frapparatus{V-FR-054-001} and which probably constitutes, at the very least,
a new step forward toward formulating an algebraic
characterization of $\mathrm{Gal}(\bar{\mathbb{Q}}/\mathbb{Q})$. This, together with the ``fundamental
conjecture'' (mentioned in par. 3 below), now appears as
the principal open question for the foundations of an ``anabelian
algebraic geometry'', which for several years has represented
(by far) my strongest focus of interest in mathematics.

\notedef{copy1-2}{(\textsuperscript{2})} I can nevertheless make an exception for another ``fact'', from the time
when, at around twelve years of age, I was interned in the concentration
camp of Rieucros (near Mende). It was there that I learned, from an
internee, Maria, who voluntarily gave me private lessons,

\margin{54'}

\finishscan
% END PAGE esq_119_054
% BEGIN PAGE esq_119_055
\scanpage{119/55}
\mspage{53}

the definition of the circle. It had impressed me by its
simplicity and obviousness, whereas the circle's property of ``perfect
roundness'' had previously seemed to me a mysterious reality
beyond words. It was at that moment, I believe, that I
glimpsed for the first time (without, of course, formulating it to myself
in these terms) the creative power of a ``good'' mathematical
definition, of a \uline{formulation} that describes the essence. Even today,
it seems that the fascination this power has exercised over
me has lost none of its force.

\notedef{copy1-3}{(\textsuperscript{3})} More generally, beyond so-called ``anabelian'' varieties
over fields of finite type, anabelian algebraic geometry
(as it emerged a few years ago) leads to a
description, in terms of profinite groups alone, of the category
of schemes of finite type over the absolute base $\mathbb{Q}$ (or
even $\mathbb{Z}$\frapparatus{V-FR-055-001}), and thereby, in principle, of the category of
arbitrary schemes (by suitable passages to the limit). It
is thus a construction that ``pretends'' to ignore rings
(such as $\mathbb{Q}$, algebras of finite type over $\mathbb{Q}$, etc.) and
the algebraic equations traditionally used to describe
schemes, by working directly with their étale topoi,
expressible in terms of systems of profinite groups. A
grain of salt, however: to be able to hope to reconstruct a
scheme (of finite type over $\mathbb{Q}$, say) from its étale topos,
which is a purely topological invariant, one should place oneself
not in the category of schemes (of finite type over $\mathbb{Q}$ in
this case), but in the one obtained from it by ``localization'',
making invertible the morphisms that are ``universal
homeomorphisms'', i.e. that are finite, radicial, and surjective. The
development of such a translation of a ``geometric world''
(namely that of schemes, scheme-theoretic multiplicities, etc.) into
terms of an ``algebraic world'' (that of profinite groups, and
systems of profinite groups describing suitable topoi (called ``étale''))
may be regarded as an ultimate outcome

\margin{55}

\finishscan
% END PAGE esq_119_055
% BEGIN PAGE esq_119_056
\scanpage{119/56}
\mspage{54}

of Galois theory, probably in the very spirit of Galois.
The perpetual question ``and what is all this for ?'' seems to me to have
neither more nor less meaning in the case of the anabelian algebraic
geometry now coming into being than for Galois theory in
Galois's time (or even today, when the question is asked
by an overwhelmed student ...), and likewise for the comment
that generally accompanies it: ``all that is very general !''.

\notedef{copy1-4}{($^4$)} It is therefore easy to conceive that a group such as $S\ell(2,\mathbb{Z})$, with
its ``arithmetic'' structure, is a veritable machine for constructing
``motivic'' representations of $\mathrm{Gal}(\bar{\mathbb{Q}}/\mathbb{Q})$ and of its
open subgroups, and that one thus obtains, at least in
principle, all motivic representations of weight 1,
or contained in a tensor product of such representations
(which already makes quite a bundle !). In 1981 I had begun
to experiment with this machine in a few particular cases,
obtaining various remarkable representations of $\mathrm{I}\!\Gamma$ in
groups $G(\mathbb{Z}^{\wedge})$, where $G$ is a group scheme (not necessarily
reductive) over $\mathbb{Z}$, starting from suitable homomorphisms
\[ S\ell(2,\mathbb{Z}) \longrightarrow G_o(\mathbb{Z}) , \]
where $G_o$ is a group scheme over $\mathbb{Z}$, and $G$ is constructed
from this as an extension of $G_o$ by a suitable group
scheme. In the ``tautological'' case $G_o = S\ell(2)_{\mathbb{Z}}$,
one finds for $G$ a remarkable extension of $G\ell(2)_{\mathbb{Z}}$ by a torus of
dimension 2, with a motivic representation that ``caps''\tnote{TN-056-01}{The quoted verb ``caps'' translates the figurative ``coiffe'', literally puts a cap on.} those
associated with the class fields of the extensions $\mathbb{Q}(i)$ and $\mathbb{Q}(j)$
(as chance would have it, the ``complex multiplication fields'' of the two
``anharmonic'' elliptic curves). There is a construction principle here
that seemed to me very general and very effective, but
I have not had (or taken) the leisure to work through its dévissage and follow it
to the end --- this is one of the many ``hot spots'' in the
program of foundations for anabelian algebraic geometry (or
``Galois theory'' in an enlarged version) that I propose to
develop. At present, and in an order of priority that is

\margin{56}

\finishscan
% END PAGE esq_119_056
% BEGIN PAGE esq_119_057
\scanpage{119/57}
\mspage{55}

probably very provisional, these points are:

a) Combinatorial construction of the Teichmüller Tower.

b) Description of the automorphism group of the profinite
compactification of this tower, and reflection on a characterization
of $\mathrm{I}\!\Gamma = \mathrm{Gal}(\bar{\mathbb{Q}}/\mathbb{Q})$ as a subgroup of the latter.

c) The ``machine for motives'' $S\ell(2,\mathbb{Z})$ and its variants.

d) The anabelian dictionary and the fundamental conjecture
(which is perhaps not so ``out of reach'' as all that !). Among
the crucial points of this dictionary, I envisage the ``profinite
paradigm'' for the fields $\mathbb{Q}$ (cf. b)), $\mathbb{R}$ and $\mathbb{C}$, for which a plausible
formulation remains to be drawn out, as well as a description of the
inertia subgroups of $\mathrm{I}\!\Gamma$, through which the passage from characteristic
zero to characteristic $p > 0$, and to the absolute ring $\mathbb{Z}$, begins.

e) Fermat's problem.

\notedef{copy1-5}{($^5$)} I shall nevertheless mention a more delicate piece of work (apart from
the work on cubic complexes mentioned in passing), on
the combinatorial interpretation of the regular maps associated with
the congruence subgroups of $S\ell(2,\mathbb{Z})$. This work was developed
above all with a view to expressing the ``arithmetic'' action of $\mathrm{I}\!\Gamma = \mathrm{Gal}(\bar{\mathbb{Q}}/\mathbb{Q})$
on these ``\uline{congruence maps}'', which takes place, essentially,
through the cyclotomic character of $\mathrm{I}\!\Gamma$. A starting
point was the combinatorial theory of the ``bi-icosahedron'' developed
in a C4 course from purely geometric motivations,
which (as it subsequently turned out) makes it possible to express conveniently
the action of $\mathrm{I}\!\Gamma$ on the category of icosahedral maps (i.e.
congruence maps of index 5).

\notedef{copy1-6}{($^6$)} Let us point out in this connection that the isomorphism classes of
compact tame spaces are \unsure{.}the same\tnote{TN-057-01}{The doubtful dot before ``les'' in ``les mêmes'' is retained before the corresponding ``the same''.} as in the ``piecewise
linear'' theory (which, I recall, is \uline{not} a tame theory).
This is, in a sense, a rehabilitation of the ``Hauptvermutung'',\tnote{TN-057-02}{The German name ``Hauptvermutung'' means ``main conjecture''. It is retained as the name used in the source.}
which is ``false'' only because, for historical reasons
that it would probably be interesting to examine more closely, the

\margin{57}

\finishscan
% END PAGE esq_119_057
% BEGIN PAGE esq_119_058
\scanpage{119/58}
\mspage{56}

foundations of topology used to formulate it did not exclude
phenomena of wildness. It goes (I hope) without saying
that the need to develop new foundations for
``geometric'' topology by no means excludes the fact that these phenomena,
like everything under the sun, have their reason for being
and their own beauty. More adequate foundations will not eliminate
these phenomena, but will allow us to place them where they
belong, as ``limiting cases'' of phenomena of ``true''
topology.

\notedef{copy1-7}{($^7$)} In fact, to reconstruct this system of spaces
\[ (i_o,\ldots,i_n) \longmapsto X_{i_o,\ldots,i_n} \]
contravariant on Drap(I) (for inclusion of flags), it
suffices to know the $X_i$ (or ``\uline{unfolded strata}'') and the $X_{i,j}$
(or ``\uline{connecting tubes}'') for $i,j \in I$, $i < j$, and the morphisms
$X_{i,j} \longrightarrow X_i$ (which are ``bordering'' inclusions)\tnote{TN-058-01}{``Bordering'' renders the author's quoted ``bordantes''; the corresponding adjective appears again with ``embedding'' on the following page.} and $X_{i,j} \longrightarrow X_i$
(which are proper fibrations, whose fibers $F_{i,j}$ are
called ``\uline{connecting fibers}'' for the strata with indices $i$ and $j$).
In the case of a tame multiplicity, however, one must
also know the ``\uline{junction spaces}'' $X_{i,j,k}$ $(i < j < k)$ and their
morphisms into $X_{i,j}$, $X_{j,k}$, and above all $X_{i,k}$, fitting into
the following commutative hexagonal diagram, in which the two squares on
the right are cartesian, the arrows $\hookrightarrow$ are immersions (not
necessarily embeddings here), and the other arrows are
proper fibrations:

\diagram{
\begin{tikzcd}[ampersand replacement=\&]
 \& X_{ik} \arrow[rr, hook] \hand{\arrow[dl]} \& \& X_k \& \\
X_i \& \& X_{ijk} \hand{\arrow[ul, hook]} \arrow[rr, hook] \hand{\arrow[dl]} \& \& X_{jk} \hand{\arrow[ul, hook]} \hand{\arrow[dl]} \\
 \& X_{ij} \hand{\arrow[ul]} \arrow[rr, hook] \& \& X_j \&
\end{tikzcd}
}

(NB. This diagram defines $X_{ijk}$ in terms of $X_{ij}$ and $X_{jk}$ over $X_j$,
but not the arrow $X_{ijk} \longrightarrow X_{ik}$, because $X_{ik} \longrightarrow X_k$ is \uline{not} necessarily
an embedding.)

\margin{58}

\finishscan
% END PAGE esq_119_058
% BEGIN PAGE esq_119_059
\scanpage{119/59}
\mspage{57}

In the case of stratified tame spaces properly so called (which
are not multiplicities strictly speaking), one can express
the ``unfolding'' of this structure conveniently, i.e.
the system of spaces $X_{i_o,\ldots,i_n}$, in terms of the tame
space $X_*$\frapparatus{V-FR-059-001} that is the sum of the $X_i$, which is endowed\frapparatus{V-FR-059-002} with the \uline{structure of an}
\uline{ordered object} (in the category of tame spaces), having as the graph
$X_{**}$ of the order relation the sum of the $X_{i,j}$ and the $X_i$
(the latter constituting the diagonal). Among the essential
properties of this ordered structure, let us note only
here that $pr_1 : X_{**} \longrightarrow X_*$ is a proper (locally trivial)
fibration, and $pr_2 : X_{**} \longrightarrow X_*$ is a ``bordering'' embedding.\frapparatus{V-FR-059-003}
There is an analogous interpretation of the unfolding of a stratified
tame multiplicity, in terms of a \uline{category structure}
(replacing a simple ordered structure) ``in the sense of tame
multiplicities'', whose composition map is given
by the morphisms $X_{ijk} \longrightarrow X_{ik}$ above.

\margin{59}

\finishscan
% END PAGE esq_119_059
\typeout{EDITION-APPARATUS-COUNT: \theappnote}
\end{document}
